\providecommand{\LegalMathRoot}{}
\providecommand{\LegalMathBibliography}{references}
\documentclass[12pt,a4paper,oneside,openany]{book}
\usepackage[margin=25mm,headheight=15pt]{geometry}
\usepackage{fontspec}
\setmainfont{TeX Gyre Pagella}
\setsansfont{TeX Gyre Heros}
\setmonofont[Scale=0.84]{DejaVu Sans Mono}
\usepackage{amsmath,amssymb,mathtools,amsthm}
\usepackage{microtype,setspace}
\setstretch{1.06}
\usepackage{booktabs,longtable,tabularx,array}
\usepackage{graphicx,xcolor,tikz,float}
\usetikzlibrary{arrows.meta,positioning,calc,shapes.geometric}
\usepackage{enumitem,listings,titlesec}
\usepackage[round,authoryear]{natbib}
\usepackage[hyphens]{url}
\Urlmuskip=0mu plus 2mu
\definecolor{navy}{HTML}{173A53}
\definecolor{teal}{HTML}{186C75}
\definecolor{pale}{HTML}{F1F5F6}
\usepackage{xr-hyper}
\usepackage[unicode,colorlinks=true,linkcolor=navy,citecolor=navy,urlcolor=navy]{hyperref}
\hypersetup{pdftitle={LegalMath: Technical Companion},pdfauthor={Chak Wong},pdfsubject={SFC circular interpretation, legal argumentation, ensemble assurance and verified Java delivery}}
\externaldocument[][nocite]{monograph}
\titleformat{\chapter}[display]{\sffamily\bfseries\color{navy}\raggedright\hyphenpenalty=10000}{\large Chapter \thechapter}{8pt}{\huge}
\titlespacing*{\chapter}{0pt}{0pt}{24pt}
\titleformat{\section}{\Large\sffamily\bfseries\color{navy}\raggedright}{\thesection}{0.7em}{}
\titleformat{\subsection}{\large\sffamily\bfseries}{\thesubsection}{0.65em}{}
\setlength{\parindent}{0pt}
\setlength{\parskip}{0.45em}
\setlength{\emergencystretch}{3em}
\setlist{itemsep=0.15em,topsep=0.35em,leftmargin=1.6em}
\setcounter{tocdepth}{1}
\makeatletter
\def\ps@legalmath{%
  \def\@oddhead{\small\sffamily\color{navy}LegalMath\enspace|\enspace Regulatory specification and execution\hfil\leftmark}%
  \let\@evenhead\@oddhead
  \def\@oddfoot{\hfil\small\thepage\hfil}%
  \let\@evenfoot\@oddfoot}
\makeatother
\renewcommand{\chaptermark}[1]{\markboth{Chapter \thechapter}{}}
\pagestyle{legalmath}
\lstset{basicstyle=\small\ttfamily,backgroundcolor=\color{pale},frame=single,rulecolor=\color{pale},breaklines=true,columns=fullflexible,keepspaces=true,showstringspaces=false,aboveskip=0.8em,belowskip=0.8em}
\newcommand{\True}{\mathsf{T}}
\newcommand{\False}{\mathsf{F}}
\newcommand{\Unknown}{\mathsf{U}}
\newcommand{\Conflict}{\mathsf{C}}
\newcommand{\OO}{\mathsf{OUT\_OF\_SCOPE}}
\newcommand{\HKD}{\mathrm{HKD}}
\providecommand{\llbracket}{\mathopen{[\![}}
\providecommand{\rrbracket}{\mathclose{]\!]}}
\newcommand{\caseheading}[1]{\subsubsection*{#1}}
\newcommand{\IR}{\mathrm{IR}}
\newcommand{\CNL}{\mathrm{CNL}}
\newcommand{\eval}{\operatorname{eval}}
\newcommand{\srcpath}[1]{\path{#1}}
\newcolumntype{Y}{>{\raggedright\arraybackslash}X}
\newcolumntype{P}[1]{>{\raggedright\arraybackslash}p{#1}}
\newtheorem{proposition}{Proposition}[chapter]
\newtheorem{definition}{Definition}[chapter]

% The same visual grammar is used for related concepts throughout the book.
\tikzset{teachbox/.style={draw=teal,fill=pale,rounded corners=2pt,
  text width=41mm,minimum height=17mm,align=center,font=\small,inner sep=4pt},
  teacharrow/.style={-{Latex[length=2mm]},draw=navy,thick}}
\newcommand{\TeachingFlow}[3]{%
\begin{tikzpicture}[node distance=8mm]
\node[teachbox] (a) {#1};
\node[teachbox,right=of a] (b) {#2};
\node[teachbox,right=of b] (c) {#3};
\draw[teacharrow] (a)--(b);\draw[teacharrow] (b)--(c);
\end{tikzpicture}}
\newcommand{\TeachingBranch}[3]{%
\begin{tikzpicture}[node distance=9mm and 11mm]
\node[teachbox] (a) {#1};
\node[teachbox,below left=of a] (b) {#2};
\node[teachbox,below right=of a] (c) {#3};
\draw[teacharrow] (a.south)--++(0,-4mm)-|(b.north);
\draw[teacharrow] (a.south)--++(0,-4mm)-|(c.north);
\end{tikzpicture}}
\newcommand{\TeachingCompare}[3]{%
\begin{tikzpicture}[node distance=8mm]
\node[teachbox] (a) {#1};
\node[teachbox,right=of a] (b) {#2};
\node[teachbox,right=of b] (c) {#3};
\draw[draw=teal,dashed] (a.south) -- ++(0,-4mm) -| (c.south);
\draw[draw=teal,dashed] (b.south)--++(0,-4mm);
\end{tikzpicture}}
\newcommand{\TeachingCycle}[3]{%
\begin{tikzpicture}[node distance=8mm]
\node[teachbox] (a) {#1};
\node[teachbox,right=of a] (b) {#2};
\node[teachbox,right=of b] (c) {#3};
\draw[teacharrow] (a)--(b);\draw[teacharrow] (b)--(c);
\draw[teacharrow] (c.south)--++(0,-8mm)-|(a.south);
\end{tikzpicture}}



\begin{document}
\raggedbottom
\frontmatter
\chapter*{Technical companion}
\addcontentsline{toc}{chapter}{Technical companion}

This reference accompanies \emph{From SFC Circulars to Reviewable Bank Controls}.
The monograph explains the decisions, mathematics and limits of the control
method. This companion begins with the implemented investigation through the
round-four checkpoint, then retains provision inventories, field definitions,
software interfaces, historical implementation records and commands for people
reproducing or auditing the work. It uses public sources and synthetic cases.
Its execution records do not establish independent legal approval or readiness
for bank production.

The two PDFs are intended to be kept together. References between them link the
reader's explanation to its reproduction record. All mathematical derivations
needed for the main argument remain in the monograph.

Chak Wong\hfill 24 September 2026

\tableofcontents
\mainmatter
\appendix
\titleformat{\chapter}[display]{\sffamily\bfseries\color{navy}\raggedright}{\large Appendix \thechapter}{8pt}{\huge}
\renewcommand{\chaptermark}[1]{\markboth{Appendix \thechapter}{}}
\chapter{Reproducing the implemented investigation}
\label{app:current-interpretation}

This chapter records the interpretation and integration capabilities accepted
through round four on 24 September 2026. It complements the explanations in
the main volume with exact entry points and retained evidence. Earlier scripted
increments are historical records in Appendix~\ref{app:full-implementation-record}.
Later implementation work has its own dated reports; the counts here belong to
this checkpoint and must not be read as a live resource balance.

\section{Which route is implemented}

The service can generate isolated readings, investigate alternatives with bounded
breadth-first or UCT search, check independently collected source claims, preserve
unsupported proposals and construct draft Java controls. The source inventories
are collected separately from candidate generation; independence of requests
does not establish statistical independence of errors. Composition combines
explicitly assigned component controls. Executed scenarios can rank questions
when their facts and outcomes are comparable. Registered official examples can
supply conditional arguments. A local scheduler revisits affected controls after
source changes.

The evidence addresses three targets: retained English and applicability,
proposed interpretation and formal expression, and generated Java behaviour.
Exact bytes and quotes support the first target's identity. Model judgments
and source comparisons inform the second. Java replay checks the third within
the exercised domain. No completion flag promotes one target into another.
The full scope statement is retained at
\path{docs/implementation/interpretation-round4/method-boundaries.md}.

Rounds two and three contain live development work. The accepted round-three
investigations retained 90 claim/candidate results for 23EC46 and 384 for 24EC16.
Of the latter, 224 are deterministic unavailable-program findings for seven
unencodable proposals. They are not model judgments about legal entailment.
The other 160 received source comparisons. All 474 pairs were accounted for,
with both investigations unresolved and ineligible for release. Additional
concerns and residual questions include programming, source and factual gaps;
their counts are not counts of independent legal ambiguities.

Round four revalidated this evidence and used controlled or retained inputs,
with no new model calls. Its final regression record contains 503 tests, zero
failures, errors or skips, and two existing dependency warnings. The source-update
exercise also ran a separate Java host. Suite counts from successive phases
overlap; summing them would misdescribe the evidence.

\section{Resolve an authority edition}

The catalog at \path{.localresources/sfc-authorities/b1/catalog.json} links an
issuing authority to editions, provision locations and supported aliases. A
citation in a candidate must resolve to a retained edition and exact text, not
merely a similar title. The catalog records publication and date evidence
separately from the requested assessment time. The accepted exercise covered
four official documents and six locator/date resolutions, including the
September 2023 and January 2026 Code editions.

The FAQ's title and its later answer updates require special care. The retained
page is not an independently archived 2010 edition. Its historical applicability
remains conditional, and remote statutory definitions remain unacquired.
The extraction checks compare two text extractions and preserve disagreement
or image warnings. OCR is not implemented. Matching extraction results cannot
establish that a scanned qualification has been read.

An assurance source manifest may name \texttt{authority\_registry},
\texttt{example\_registry} and \texttt{example\_scope}; relative paths resolve
against that manifest. The scope states issuer, jurisdiction and issue IDs.
The public example
\path{examples/interpretation-assurance/authority-and-examples.json}
uses the retained catalog and example registry. It is a development input,
not an institution's approved source register.

\section{Compose separate controls}

The component specification binds the source packet, inventoried claims and
complete final candidate set. Each component states its selected claims,
alternative candidate IDs, output meaning and assignment basis. The caller
must justify why the members are rival readings of one control, or separate
controls that belong together. The tool does not infer that legal grouping
from successful compilation.

The command's required arguments are:
\begin{lstlisting}
.venv/bin/python -m legalmath.cli assurance-compose \
  --investigation INVESTIGATION_DIRECTORY \
  --spec COMPONENT_SPECIFICATION.json \
  --jdk JDK_DIRECTORY --out NEW_OUTPUT_DIRECTORY
\end{lstlisting}
These capitalized arguments stand for operator-supplied paths. Use a fresh
output directory. A retained real-input specification is
\path{artifacts/interpretation/round4/B2a/attempt-01/composition/23ec46-spec.json};
its identities belong to the referenced historical investigation and cannot
be reused after changing the source or candidate.

The output includes bundles, generated Java, JARs, cases, conformance results
and a manifest. Alternative selections form separate packages. Within each
package, rules and facts use component namespaces so unrelated meanings cannot
collide. Each result retains its source claims, modality, scope, exceptions,
timing and upstream uncertainty. Missing claims, unsupported expressions,
unassigned candidates and combinations deferred by the work limit remain
explicitly incomplete. The package is a draft.

Seven builds were exercised: two controlled two-component selections and five
retained gift-control alternatives. True, false, unknown, conflict and
out-of-scope results were covered. These are not seven whole-circular
implementations; complete coverage of 23EC46 and 24EC16 remains unproved.

\section{Run a case that distinguishes readings}

The question procedure begins with a source-bound candidate comparison and its
factual witness. The record includes the source packet, candidate set, assessment
time, base candidate, subject, output convention, assumptions and source quotes.
Before replay, the implementation checks subject and output compatibility and
constructs the declared fact correspondence. The original candidate programs
then run in Python and Java on the corresponding snapshots.

The partition record retains accepted replays, excluded candidates and reasons,
work consumed and separated pairs. Unit-weight scores count executed pairs
with different projected results. They are scheduling observations, not legal
probabilities. Resource-deferred work cannot be labelled complete; unsupported
programs and incompatible facts cannot earn a score. Source-inventory omission
checks retain priority even if every comparable output agrees.

The controlled gift/discount case separated one pair and reached the repair
scheduler. The retained 23EC46 comparison had incompatible fact bindings;
24EC16 had no checked pair. Both reported \texttt{NO\_CHECKED\_DISTINCTIONS}.
This means no checked distinction was available, not that all possible readings
agreed. Conditional derived comparisons remain a separate facility. The
implemented partition and verifier are in
\path{src/legalmath/interpretation/assurance/questions.py}; retained execution
is indexed by \path{docs/implementation/interpretation-round4/b2b-result.md}.

\section{Use an official example without inventing authority}

The register at \path{.localresources/sfc-examples/b3/registry.json} binds each
example to the authority catalog, exact source quotation, issuer, edition,
proposed factors and declared outcome. The current issue is
\texttt{paragraph.3.11.applicability}. It concerns whether the paragraph applies,
not an unconditional prohibition after exceptions.

Two proposed branches of FAQ Question 1 produced twenty conditional analogy
and distinction moves on five candidates. They share one source family. Six
questions remained, including historical applicability and whether the selected
factors were legally material. The factorization is a developer proposal that
requires review; genuine quoted words do not make it exhaustive.

A generated example outside the register remains an unregistered hypothetical
proposal. A quotation cannot authenticate an invented court, case identity or
holding. This implementation has no court-case corpus, general judicial-ratio
extractor or procedure for determining binding priority. Source or registry
changes invalidate dependent analysis instead of silently retaining its former
applicability.

\section{Present proposals and preserve partial checks}

Present candidate analysis through \texttt{candidate-presentation.json}.
An entry with \texttt{WITHHELD\_CORRUPT\_PROSE} has a null presentation and points
to \texttt{prose-quality.json} and the original candidate. Raw proposals,
compiled programs and their source checks remain retained. An entry with
\texttt{NO\_PATTERN\_DETECTED} passes only the implemented pattern diagnostic.
It has no certificate of natural language coherence or legal correctness.

The detector identified five defects in the actual retained 24EC16 proposal:
placeholder subject and statement fields, self-repair chatter, and output
instructions embedded in questions. Exact authentic quotations, ordinary
uncertainty and non-English text are preserved by the checked negative cases.
The detector's coverage is deliberately narrow. Editing the proposal would
require a successor candidate and renewed checks; the presentation mechanism
does not silently rewrite its meaning.

Fidelity results are checked for exact claim/candidate pairs, source quotations
and representation quotations. The historical 23EC46 response contained 33 rows
where 32 were required; the validator identified the extra pair, and the
retained correction validated all 32. That was revalidation of an existing
correction, not a new live repair.

The per-response schema permits at most 1,000 checks and 30 concerns. The
aggregate profile bounds an investigation at 4,096 pairs, 1,080 concerns and
2,097,152 bytes; its per-call pair default is 32 with a maximum of 64. These
are engineering capacity choices. Admission considers cached and unsupported
readings as well as pending model work. It can conservatively refuse work that
might have returned shorter responses.

A completed matrix is retained in \texttt{complete.json}. Interrupted aggregation
keeps \texttt{partial.json} and the validated parts; an oversized cached part
remains separately available. Repeated concerns are retained even when a
separate presentation list deduplicates residual questions. A later failure
cannot erase earlier rows or turn a partial result into a clean interpretation.
The full profile and record names are in
\path{docs/implementation/interpretation-round4/operator-guide.md}.

\section{Freeze and execute a comparison}

The implemented adapters are \texttt{single-reader}, \texttt{isolated-readers},
\texttt{search} and \texttt{assurance}. Freeze the source, assessment time,
case-family assignments, reference outcomes and resource ceilings through
\path{legalmath.interpretation.assurance.evaluation.freeze} before dispatch.
The hidden scoring references are unavailable to the methods. Missing answers
and abstention are reported alongside unsafe and useful accepted answers.
The unit of uncertainty is a paired source-family average, as derived in
Section~\ref{sec:family-comparison} of the monograph.

A replay invocation has this form:
\begin{lstlisting}
.venv/bin/python -m legalmath.cli assurance-evaluate \
  --frozen FROZEN_STUDY.json --out NEW_OUTPUT_DIRECTORY \
  --jdk JDK_DIRECTORY --replay-responses RESPONSES.json
\end{lstlisting}
For authorized live work, the provider arguments are instead
\texttt{--allowance LEDGER.json --total-calls AUTHORIZED\_CEILING}.
The command checks the complete reservation before dispatch. Supplying the
flag does not increase the authorized ledger ceiling.

The round-four fixture reserved eight calls per method for two cases, or 64
calls. At that checkpoint only 22 of the original 100 remained. Consequently
no live comparison was run. The eight scripted jobs tested actual adapters,
hidden finite-scenario scoring and adverse outcomes. Their one development
family supplied no informative statistical separation and no estimate of
English error or reviewer savings. Additional calls alone would not make
those cases a held-out sample.

Subsequent work on references and source families is recorded separately in
\path{docs/implementation/interpretation-round5/execution-report.md}. Those reference preparations and new-source attempts form a separate
continuation. Historical call balances above must not be copied into a new run; use
the actual ledger and that run's reviewed study. A complete effectiveness study
still needs defensible outcome evidence and a justified sampling and precision
design. Internal qualified reviewers or explicit official outcomes can support
suitable tasks; there is no automatic external-consultancy requirement.

\section{Package Java and revisit source changes}

The package builder retains the draft rule bundle, policy source, JAR, host
caller, source commitments, conformance cases and manifest. A separate Java
process calls the generated public API and receives the full typed result.
Unknown or conflict must not be reduced to a Boolean authorization. The local
exercise uses draft mode and supplies no institution release identity.

The scheduler wraps the fixed assurance monitor with bounded visits and process
timeouts. Its command accepts:
\begin{lstlisting}
.venv/bin/python -m legalmath.cli assurance-schedule \
  --config REGISTRY.json --resource-root RESOURCE_DIRECTORY \
  --out NEW_OUTPUT_DIRECTORY --jdk JDK_DIRECTORY \
  --at ASSESSMENT_TIME --settings SETTINGS.json \
  --replay-responses RESPONSES.json \
  --ticks 1 --interval-seconds 60 --tick-timeout-seconds 600
\end{lstlisting}
The numbers shown illustrate a bounded replay invocation, not a recommended
bank cadence. The default is one visit. The assessment time stays fixed during
an invocation; the institution scheduler must supply the intended assessment
time for subsequent invocations. Authorized live provider arguments use the
same allowance and total-call pair as evaluation.

Before dispatch, the wrapper copies source, registry, settings and replay bytes
into the output record and binds their identities. Editing an input file while
a visit runs cannot change its frozen inputs. Failed or uncertain work consumes
an attempt and remains visible. Nested failure or unresolved meaning produces
an attention state even when the outer process exits successfully.

The final integration repair prevents an unperformed example query, a
resource-deferred partition or limited argument analysis from being cached as
fully executed. Completed but unresolved interpretation remains unresolved on
unchanged visits. Retries are bounded, and corrupted cached evidence is rejected.
The status distinguishes what ran from what was established.

The accepted update exercise launched three separate Java host processes.
Adding a fee-discount exception changed the selected restriction from true to
false on the controlled facts. An unaffected control and the earlier package
were preserved. Corrupted packages, uncertain callbacks, missed cadence and
retry exhaustion were also exercised. No actual regulatory amendment was
independently interpreted in this demonstration.

For institution staging, the bank must supply its host and JDK, service identity,
control register and fact bindings, operations-owned scheduler, update cadence
and release roles. The local command installs no service and activates no bank
policy. Confidentiality, authenticated release, audit integrity and concurrent
host transactions require evidence from that environment.

\clearpage
\section{Find the retained evidence}

The acceptance state is in
\path{artifacts/interpretation/round4/state.json}. Each accepted phase points to
an immutable attempt directory and its \texttt{run-manifest.json}; inspect the
manifest for exact commands, inputs, logs and output identities. The compact
index below explains what each record establishes.

\begin{description}[style=nextline]
\item[B0: proposal and matrix integrity]
Actual corrupt prose and the extra historical response pair, plus capacity and
partial-result retention. See \path{docs/implementation/interpretation-round4/execution-report.md}.
\item[B1: source editions]
Four documents and six locator/date resolutions, with missing definitions and
historical uncertainty retained. See \path{docs/implementation/interpretation-round4/b1-result.md}.
\item[B2a and B2b: composition and distinguishing cases]
Seven Java builds and the controlled executed partition, with no comparable
witness claimed for either retained circular. See the corresponding result
notes in \path{docs/implementation/interpretation-round4/}.
\item[B3 and B4: examples and evaluation]
One official-example family and eight scripted comparison jobs. These establish
conditional argument construction and evaluation mechanics. Their result notes
state the remaining semantic and empirical limits.
\item[B5 and B6: integration and completion]
Local Java-host update and final status/input-retention repairs. The final suite
is \path{artifacts/interpretation/round4/B6/attempt-01/tests.xml}; its adjacent
manifest binds the 503-test acceptance and operations exercise.
\end{description}

The preceding live development is documented in the round-two and round-three
execution reports under \path{docs/implementation/}. These records explain how
proposals and source judgments were obtained; round-four replay does not turn
them into fresh independent observations. The two retained circulars remain
unresolved. No method ranking, population interpretation accuracy or reduction
in review cost is established by this evidence index.

% Detailed tables retained for implementers and auditors. The main narrative
% introduces their questions and uses smaller reader-facing summaries first.
\chapter{Technical implementation record}
\label{app:implementation-record}

The main chapters use a small number of examples to teach the decisions. This
appendix preserves the exhaustive mappings that an implementation or audit may
need to look up. The tables are reference material; they are not intended to
carry the argument without the surrounding explanations in Chapter~2.

\section{Complete SPI provision inventory}
\begingroup\small
\begin{longtable}{P{0.14\textwidth}P{0.37\textwidth}P{0.39\textwidth}}
\caption{Disposition of the complete circular for the worked profile. Subparagraphs
within each stated range remain part of that row.}\label{tab:circular-inventory}\\
\toprule Source & Substantive instruction & Treatment in the worked change \\\midrule
\endfirsthead
\toprule Source & Substantive instruction & Treatment in the worked change \\\midrule
\endhead
Main text; Annex 1 opening and footnote 1 & Introduce streamlining, retain conduct and risk controls, and import definitions and suitability requirements. & Record authority, perimeter and dependencies. The Java result covers selected conditions only; the bank retains its wider control chain.\\
1--2 & Individual criteria at the relevant date; separate investment-holding corporation route. & Individual PI assessment is an input. A corporate client returns outside profile; ownership and qualifying-owner aggregation require a separate rule family.\\
3.1(a)--(b) & Inclusive alternative wealth thresholds and currency equivalents. & Execute exact HK-cent comparisons joined by OR. Foreign currency requires an adopted conversion mapping before evaluation.\\
3.2(a)--(d), 3.3, 3.4(a)(i)--(iii), (b) & Portfolio and asset/liability attribution; residence exclusion; written or equal joint shares under the stated circumstances. & Reviewed normalization produces portfolio and net assets. A missing agreement is unresolved evidence, not established absence. Preserve each ownership calculation.\\
4.1(a)--(d), 5.1 & Reasonable sophistication assessment, alternative supporting routes and non-conservative objectives. & Execute the alternative routes and require the dated assessment and objectives. Do not substitute the representative's attributes.\\
6, 7.1--7.3 & Restrict transactions to selected categories and threshold; define categories, provide information, and apply category qualification. & Require selected category, resolved category explanations and information. Count same-category transactions or establish another qualification route.\\
8.1--8.2 & Client specifies a threshold appropriate to circumstances; bank retains the setting and rationale. & Accept an agreed, reviewed absolute HKD threshold and rationale. A percentage-of-AUM arrangement needs a versioned conversion to the same input.\\
8.3(a) & Monitor gross exposure on execution. & Compare projected exposure including the candidate transaction with the agreed limit; host must reserve exposure atomically before using the decision.\\
8.3(b) & Alternative monitoring for designated accounts and new funds. & Explicitly outside profile. Implement its own account, transferred-position and top-up events before enabling it.\\
% BEGIN REVISION ADDITION gap-spi-dispositions
\bottomrule
\end{longtable}
\begingroup\normalsize
\begin{figure}[H]
\centering
\TeachingCompare{Qualification and category}{Exposure and transaction relief}{Records, consent and continuing review}
\caption{The provision inventory preserves responsibilities beyond the financial test.}\label{fig:gap-spi-dispositions}
\end{figure}

\endgroup
\begin{longtable}{P{0.14\textwidth}P{0.37\textwidth}P{0.39\textwidth}}

\toprule Source & Substantive instruction & Treatment in the worked change \\\midrule
\endhead
% END REVISION ADDITION gap-spi-dispositions
8.4--8.5 & Detect outsize or material transactions and warn; review threshold compliance annually and alert if exceeded. & Require transaction assessment and conditional warning. Schedule review/alert tasks; a reviewed composite reports annual controls current.\\
9; 10.1--10.4 & Define solicited-transaction relief, with explanation exceptions and offering documents retained. & Execute the solicited profile. The result identifies applicable relief; it does not disable every suitability or disclosure control.\\
11.1--11.5; footnotes 2--4 & Different relief for unsolicited complex-product transactions, including document alternatives and warning arrangements. & Outside profile. Preserve the distinct document/evidence branches, bond information alternatives and annual-warning requirement for a later implementation.\\
12.1--12.3 & Reasonable satisfaction, written assessment and correspondence recording category and threshold choices. & Dated assessments and document references feed named facts; records remain independently inspectable.\\
13.1(a)--(d), 13.2 & Prior written agreement and explanation of consequences, categories, threshold and withdrawal/amendment rights. & A checklist-backed acknowledgment input plus a replayed consent state. Every item in the checklist must be evidenced; a signature alone is insufficient.\\
14.1--14.2 & Annual continued-qualification/agreement review and written reminders covering consequences, categories, threshold, breaches and rights. & Ongoing tasks and reviewed status. Suspending streamlining when review is overdue is a proposed additional bank restriction.\\
Annex 2, FAQ 1 & Appropriate reliance on disclosures after reasonable KYC efforts; clarify inconsistency. & Define admissible evidence and escalation; do not invent a blanket requirement for independent third-party confirmation.\\
FAQ 2 & Client and power-of-attorney distinctions. & Attribute qualifications to the client. Apply the FAQ to transaction-history evidence without copying the attorney's sophistication.\\
FAQs 3--4, including notes & Category design, difficult product distinctions and category information. & Bank-owned category taxonomy, versioned mappings and information checklist; uncertain mapping requests review.\\
FAQ 5 & Threshold should reflect the client's circumstances; bank-held AUM may be a limited part of wealth. & Recorded client circumstances and threshold rationale. No universal percentage is invented.\\
FAQ 6 & Explain both monitoring approaches, including designated-account transfers, appreciation and new funds. & Use execution monitoring here. Keep alternative-account rules visible; do not infer forced unwinding merely from exceeding its threshold.\\
FAQ 7 & Gross exposure includes leverage. & Feed leverage-inclusive projected exposure to the execution comparison.\\
\bottomrule
\end{longtable}
\endgroup

\section{Detailed inputs for the demonstrated Java profile}
\begingroup\small
\begin{longtable}{P{0.30\textwidth}P{0.60\textwidth}}
\caption{Input contracts for the demonstrated Java profile. A known value always
has evidence, a validity interval and a recording time.}\label{tab:demo-inputs}\\
\toprule Input or group & Meaning and intended owner \\\midrule
\endfirsthead
\toprule Input or group & Meaning and intended owner \\\midrule
\endhead
Profile flags & \texttt{profile\_individual}, \texttt{profile\_solicited}, \texttt{profile\_execution\_monitoring}: the host's classified client, activity and monitoring mode. False selects an unsupported profile; unknown requests review.\\
\texttt{individual\_pi} & Compliance assessment under the adopted Professional Investor Rules and Code definitions, for this legal client and date. It is distinct from merely being an individual.\\
\texttt{portfolio}; \texttt{net\_assets\_ex\_home} & Valuation/operations supplies exact HK cents after reviewed ownership, liabilities, residence and currency treatment. Preserve the calculation inputs and source evidence.\\
Qualification fields & \texttt{relevant\_degree}, \texttt{professional\_qualification}, \texttt{relevant\_year\_of\_work}; alternatively \texttt{category\_transactions\_3y}. Client-level evidence is mapped to the selected category and relevant period.\\
Assessment fields & \texttt{reasonable\_sophistication\_assessment}, \texttt{non\_conservative\_objectives}, \texttt{written\_assessment\_retained}. Compliance/relationship management supplies the assessment and records; missing review is unknown.\\
Category fields & \texttt{category\_selected}, \texttt{category\_information\_delivered}, \texttt{category\_explanation\_resolved}. Client choice and delivery evidence must match the transaction's versioned category.\\
% BEGIN REVISION ADDITION gap-spi-inputs
\bottomrule
\end{longtable}
\begingroup\normalsize
\begin{figure}[H]
\centering
\TeachingCompare{Reviewed client attributes}{Transaction and exposure evidence}{Consent and completed actions}
\caption{The input table joins facts supplied by different operational owners.}\label{fig:gap-spi-inputs}
\end{figure}

\endgroup
\begin{longtable}{P{0.30\textwidth}P{0.60\textwidth}}

\toprule Input or group & Meaning and intended owner \\\midrule
\endhead
% END REVISION ADDITION gap-spi-inputs
Acknowledgment and consent & \texttt{prior\_written\_acknowledgment\_complete} covers every required item in 13.1--13.2; \texttt{active\_consent} comes from the client/category consent history. Document completeness and present consent are different facts.\\
Threshold and exposure & \texttt{streamlining\_threshold}, \texttt{projected\_gross\_exposure} are HK cents; \texttt{threshold\_rationale\_retained} is an assessment/record check. Order management supplies a consistent exposure snapshot including leverage.\\
\texttt{annual\_review\_current} & Reviewed composite covering continued qualification/agreement, threshold review, required reminders and breach alerts. The demo consumes this status; it does not calculate annual deadlines or send reminders.\\
Product information & \texttt{current\_offering\_documents\_delivered}, \texttt{explanation\_requested\_or\_material\_query}, \texttt{explanation\_delivered}. Delivery/interaction records determine whether the remaining request has been answered.\\
Transaction warnings & \texttt{transaction\_materiality\_assessed}, \texttt{outsize\_or\_material\_transaction}, \texttt{material\_warning\_delivered}. The assessment is required; a warning is required if the assessment says so.\\
\bottomrule
\end{longtable}
\endgroup

\section{Full distinguishing-case table}
\begingroup\small
\begin{longtable}{P{0.38\textwidth}P{0.19\textwidth}P{0.33\textwidth}}
\caption{Distinguishing cases derived from the running interpretation. Other
inputs remain as in the successful synthetic case unless stated.}\label{tab:dry-cases}\\
\toprule Changed fact & Result & Why it changes the operational response \\\midrule
\endfirsthead
\toprule Changed fact & Result & Why it changes the operational response \\\midrule
\endhead
Portfolio HK\$39,999,999.99; net assets HK\$75m & Standard process & Both financial alternatives are refuted.\\
Portfolio HK\$35m; net assets HK\$90m & Conditions met & The alternative net-asset route suffices.\\
Portfolio missing; net assets HK\$90m & Conditions met & Missing portfolio is reported, but does not block the established alternative.\\
Portfolio missing; net assets HK\$70m & Review required & Portfolio evidence can change the answer. It is a blocking input.\\
Conservative objectives & Standard process & Wealth does not overcome paragraph 5.1.\\
Four category transactions; no other qualifying route & Standard process & The countable route is not established.\\
Chosen category does not match & Standard process & General SPI attributes do not extend the selected categories.\\
Gross exposure HK\$10,000,000.01 & Standard process & The execution-monitoring condition is exceeded.\\
Written acknowledgment incomplete, or consent withdrawn & Standard process & Prior documentation and current consent are both required.\\
Client requests explanation; no answer recorded & Standard process & Paragraph 10.3 retains the requested explanation.\\
Material transaction; no warning delivered & Standard process & The applicable paragraph 8.4 warning remains outstanding.\\
Two inconsistent portfolio records & Review required & An evidence conflict is not resolved by choosing the more convenient amount.\\
Designated-account monitoring selected & Outside profile & Its different paragraph 8.3(b)/FAQ 6 logic has not been executed by this library.\\
\bottomrule
\end{longtable}
\endgroup

% The complete engineering record retained outside the reader-facing route.
\chapter{Historical implementation and design record}
\label{app:full-implementation-record}

This appendix preserves the original workbench and the 23 September scripted
interpretation increment, together with their design proposals. Its task names,
API sketches and future-tense passages belong to those checkpoints. For the
implemented source generation, search, assurance, composition and monitoring
through 24 September, begin with Appendix~\ref{app:current-interpretation}.
Empirical effectiveness and institution acceptance remain separate requirements.


% BEGIN SOURCE UNIT P-03-product:00
\section{The shared workbench and its users}
\label{merge:P-03-product:00}

\label{sec:product}

LegalMath should be a regulatory specification workbench. Its primary users are
the compliance officer interpreting a change, the product or operations owner
deciding how it affects the business, and the engineer implementing the resulting
control. An auditor or independent reviewer should be able to reconstruct the
same decision later. The product succeeds when those people can inspect one
maintained interpretation through views suited to their work.

The first local release was designed around a narrow end-to-end task: take a selected public
circular and its dependencies, prepare proposed rules, resolve or retain
interpretation questions, evaluate synthetic cases, and export a versioned change
package containing generated Java source, a compiled JAR and its evidence.
That local path is now implemented with public sources and synthetic
identities. Integration into bank systems still follows the bank's ordinary
change process.
% BEGIN REVISION ADDITION teach-P-03-product-00
\begin{figure}[H]
\centering
\TeachingCompare{Compliance interprets}{Engineering implements}{Operations acts and retains evidence}
\caption{The workbench gives different participants views of the same reviewed meaning.}\label{fig:teach-P-03-product-00}
\end{figure}

% END REVISION ADDITION teach-P-03-product-00
% END SOURCE UNIT P-03-product:00

% BEGIN SOURCE UNIT P-03-product:01
\section{What a reviewer should see}
\label{merge:P-03-product:01}


Opening the SPI financial condition should display the source passage and its
annex context beside a controlled explanation: the portfolio route or the
net-asset route can satisfy this particular condition. Definitions should be
available in place, especially ownership attribution and exclusion of the primary
residence. The same screen should show the boundary cases, missing evidence and
where the condition is used in the larger SPI assessment.

The reviewer can change an interpretation, ask for a distinguishing example, or
record a question for the policy owner. If two candidate readings differ only at
the threshold, the workbench should produce that boundary case. If they differ
because one applies to a representative and another to the client, the example
should name the roles. The user should not need to understand a solver formula
to see why the choice matters.

The engineer views the corresponding types, data mappings, executable dependency
graph and tests. The operations owner views required actions, responsible teams,
timing and evidence of completion. A diagram or decision table is generated from
the same adopted rules. Maintaining a separate hand-edited flowchart would create
another translation boundary and another version to reconcile.

\begin{figure}[htbp]
\centering
\begin{tikzpicture}[node distance=8mm and 8mm,
 box/.style={draw=teal,fill=pale,rounded corners,text width=40mm,align=center,minimum height=17mm,font=\small},
 arr/.style={-{Latex[length=2mm]},thick,draw=navy}]
\node[box] (source) {Preserved sources\\and dependencies};
\node[box,right=of source] (draft) {Proposed interpretation\\and unresolved questions};
\node[box,right=of draft] (spec) {Reviewed typed rules\\and examples};
\node[box,below=of spec] (exec) {Reference evaluator\\and generated Java};
\node[box,left=of exec] (evidence) {Tests, counterexamples\\and scoped proofs};
\node[box,left=of evidence] (history) {Decision explanations\\and amendment history};
\draw[arr] (source)--(draft); \draw[arr] (draft)--(spec);
\draw[arr] (spec)--(exec); \draw[arr] (exec)--(evidence);
\draw[arr] (evidence)--(history);
\draw[arr] (evidence.north)--(draft.south);
\end{tikzpicture}
\caption{The proposed workbench. Counterexamples return to interpretation review.
A proof result remains attached to the exact rules, assumptions and software
version it checks.}
\label{fig:architecture}
\end{figure}
% END SOURCE UNIT P-03-product:01

% BEGIN SOURCE UNIT P-03-product:07
\section{The execution architecture and its deliberate limits}
\label{merge:P-03-product:07}


The authoring runtime is a Python 3.11 reference interpreter, with JSON Schema and
semantic validation, SQLite storage, content-addressed source files and a local
FastAPI service. A server-rendered interface presents source, rule and case
together. The interpreter walks a validated syntax tree; it does not execute
model-generated Python. Sets, quantifiers, general recursion, arbitrary code,
currency conversion and unbounded temporal formulas are outside RuleIR 0.1.
An external, evidenced normalizer can supply a reviewed count or HKD value.

These restrictions keep common questions tractable and make the trusted code
small enough to inspect. They also provide a natural place to evaluate the
literature: Catala for the decision fragment, DMN as a possible Java integration comparison,
and a Stipula-inspired state model for a bounded workflow. The project should not
integrate all candidate engines before demonstrating one complete use case.

% BEGIN REVISION ADDITION teach-P-03-product-07
\begin{figure}[H]
\centering
\TeachingFlow{Validate the restricted expression}{Run its defined semantics}{Reject unsupported constructs}
\caption{A bounded language makes its supported behavior inspectable.}\label{fig:teach-P-03-product-07}
\end{figure}

% END REVISION ADDITION teach-P-03-product-07
A reference interpreter implements the written execution semantics. A schema
validator checks structure; a semantic validator checks expression types,
resolved dependencies,
source links and explicit exception relationships; a scenario runner evaluates
reviewed examples. A solver adapter then searches for conflicts, unreachable
branches and semantic differences. Generated code should come from a deterministic
compiler for the supported fragment. At transaction time, evaluation uses an
approved rule version and a recorded fact snapshot; language-model generation
belongs to authoring and review.

An implementation adapter is accepted only after its treatment of missing
values, numeric boundaries, priority, dates and obligations agrees with the
reference semantics on the required checks. DMN's hit-policy choices and Rego's
undefined results make this comparison substantive
\citep{dmn2024,opadocs}. A backend that cannot express a required distinction
should report that limitation instead of approximating it silently.
% END SOURCE UNIT P-03-product:07

% BEGIN SOURCE UNIT P-03c-storage-review:00
\subsection{Source identity, review and storage}
\label{merge:P-03c-storage-review:00}


The source store preserves original bytes and an exact text derivative. The raw
document and derivative each have a SHA-256 digest. A source span records the
derivative hash, page, half-open character offsets and a hash of the quoted text.
Extraction changes create a new derivative; they do not silently move an approved
span. A reviewer can open the original page beside the extracted passage.
Scanned or poorly extracted clauses require visual confirmation before release.

The source register distinguishes circulars, annexes, FAQs, codes, bank policy
and interpretation. Dependency rows record references, definitions, amendments
and replacements; unresolved targets remain explicit. Publication date does not
automatically establish effectiveness. Applicability records identify legal
entity, regulated role, activity, product, client class and jurisdiction. A
product-provider requirement cannot enter the bank's distributor rules merely
because its title mentions a product the bank sells.

% BEGIN REVISION ADDITION teach-P-03c-storage-review-00
\begin{figure}[H]
\centering
\TeachingFlow{Retain original source bytes}{Bind interpretation to exact spans}{Approve a particular bundle hash}
\caption{Source identity and review identity must remain connected.}\label{fig:teach-P-03c-storage-review-00}
\end{figure}

% END REVISION ADDITION teach-P-03c-storage-review-00
The bundle contains definitions, expressions, spans and interpretations. Its
canonical JSON hash identifies the entire immutable revision. Review decisions
live outside those bytes and approve that exact hash. The lifecycle is draft,
in review, approved, released and retired; rejection returns a revision to draft,
while an edit creates a new draft hash. Release requires distinct compliance and
engineering reviewers, resolved blocking questions and dependencies, a supported
effective interval and passing evidence for that hash. The author cannot approve
their own revision under the proposed prototype policy.

SQLite transactions enforce review revisions and idempotency. A stale review
update is rejected rather than overwriting another review. Content-addressed
blobs are written before their database references; a missing referenced blob
is an integrity error. Audit rows record accepted transitions. Hashes detect
accidental corruption but do not make a database tamper-proof against an
administrator able to rewrite both records and hashes. Enterprise identity,
signed retention and access controls belong to the later bank deployment.
% END SOURCE UNIT P-03c-storage-review:00

% BEGIN SOURCE UNIT P-04-prototype:02
\section{The original work packages and their execution}

The original W packages specified the route from a document-only contract to
a working system. Their requirements are retained because they explain what the
implementation must preserve. The next section records their completed T-task
realization and identifies the optional research and human-evaluation work that
remains. The twelve-week allocation was a planning assumption, not measured
development effort or evidence of bank readiness.


\label{merge:P-04-prototype:02}


The implementation sequence below separates the delivered Java demonstration
from the complete product modules. Its operational entry point is
\path{docs/implementation/START-HERE.md}, with corresponding source paths and
commands in \path{work-packages.md}. The manuscript supplies the meanings,
interfaces and acceptance conditions; these files are the implementation copies.
W00 establishes an isolated
Python 3.11 environment and locks actual dependency versions. It does not modify
the environments of the three adjacent MCP projects. The proposed stack is
Python's exact integer/date/hash/SQLite facilities, JSON Schema, Pydantic/FastAPI,
Poppler or a ResearchAssistant extraction adapter, Jinja2 for the review pages,
and optional Z3 for authoring. Java 17, a deterministic source emitter and a
compiled conformance harness are required for delivery. Package versions are resolved and recorded during bootstrap;
the current local MVP has integrated and locked its selected dependencies.
The retained W-package descriptions below explain the original work division.

The following command works on the delivered specification pack now:
\begin{lstlisting}
python3 scripts/check_spec_pack.py
\end{lstlisting}
It validates schemas, fixture structure, selected static constraints, the source
anchor and the SQLite definition. After W03 implements the evaluator, the same
checker has an explicit runtime mode:
\begin{lstlisting}
.venv/bin/python scripts/check_spec_pack.py \
  --runtime legalmath.conformance:evaluate_case
\end{lstlisting}
The second command became an executed check in the accepted MVP: all 35
decision cases passed. Its event-and-release flag describes only this particular
harness; separate tests and the complete walkthrough execute those components.
The narrower SPI-Demo1 command in Section~\ref{sec:java} remains a separate
reproduction route. Neither result establishes automatic legal translation.

% BEGIN REVISION ADDITION teach-P-04-prototype-02
\begin{figure}[H]
\centering
\TeachingFlow{Reference semantics}{Java and review workflow}{Evaluation and integration}
\caption{The original work packages explain the order of engineering dependencies.}\label{fig:teach-P-04-prototype-02}
\end{figure}

% END REVISION ADDITION teach-P-04-prototype-02
\begingroup\small
\begin{longtable}{P{0.12\textwidth}P{0.42\textwidth}P{0.36\textwidth}}
\caption{Implementation packages and their completion evidence. Detailed files,
functions, commands and repair conditions are in the accompanying backlog.}\label{tab:phases}\\
\toprule Package & Work and demonstrable result & Evidence needed to proceed \\\midrule
\endfirsthead
\toprule Package & Work and demonstrable result & Evidence needed to proceed \\\midrule
\endhead
W00 & Package, locked environment, canonical JSON, immutable blobs and database migration. & Same content hashes across processes; duplicate keys/floats rejected; interrupted writes leave no broken references.\\
W01 & Import 23EC35 and annexes; preserve derivatives, spans and dependencies. & Reproduce the paragraph 3.1 anchor; changed bytes and broken attachments are detected.\\
W02 & Typed RuleIR, references, graph validation and dated evidence normalization. & Reject malformed/type-invalid/cyclic rules; explain missing, stale and conflicting facts.\\
W03 & Pure evaluator and deterministic trace. First financial-condition demonstration. & Execute 35 decision fixtures; catch boundary, OR, unknown and conflict mutations.\\
W03J & Complete Java emitter, runtime support, compiled conformance and host boundary. & Full supported results agree with the reference; mutations fail; source and JAR builds reproduce; a separate caller runs.\\
W04 & Exact-hash reviews, release guards, Java package export and amendment comparison. & Self-approval, stale reviews and wrong-bundle or wrong-JAR evidence fail; historic versions replay.\\
W05 & Consent, one duty profile, calendar policies and complete-window evidence. & Eight event fixtures plus ordering, duplicate, leap-date and late-correction tests.\\
W06 & Local API and source/rule/case/amendment screens. & Contract tests and a reviewer walkthrough; unresolved evidence remains visible.\\
W07 & Restricted SMT encoding and counterexample service. & Boundary counterexample replays; empty domain and unsupported constructs cannot pass as equivalence.\\
W08 & Paginated corpus tracking, clause coverage and bounded model drafting. & Drafting cannot approve; omissions scored against an independent source inventory.\\
W09 & Optional isolated ResearchAssistant, MathDevMCP and DynareMCP adapters. & Exact command/version/hash evidence; structural checks never mislabeled as proofs.\\
W10 & Registered comparative pilot, error analysis, cost and review-time study. & Source parity, independent adjudication, uncertainty and a scoped investment decision.\\
\bottomrule
\end{longtable}
\endgroup

The reference interpreter comes before language-model drafting because it supplies
a stable target for evaluating the draft. That ordering also measures whether the
shared specification is useful even without generative assistance. The event
model is introduced only after the decision language is coherent, so difficulties
with consent histories can be separated from errors in money comparisons or
source extraction.

W00--W04, including W03J, receive a provisional four-week allocation; W05--W08
receive another four, with review, integration, rework and the pilot bringing
the allowance to twelve weeks. The earlier three-week first-stage estimate did
not include a required Java backend. This revised allocation remains a planning
assumption, with scope reduction or a schedule extension if full operator
conformance takes longer. Each package specifies files, function signatures, dependent
packages and tests. Its acceptance condition concerns behavior, not elapsed time.
For example, W03 supplies \texttt{ir/evaluate.py} and \texttt{ir/trace.py}; W04
supplies \texttt{review/lifecycle.py} and \texttt{review/amendment.py}; W07
supplies \texttt{analysis/z3\_encode.py} and \texttt{analysis/compare.py}.

The deployment interface is the generated Java library, checked against the
reference decision API. Catala, L4 and DMN
are later adapter comparisons against the same fixed cases, not prerequisites
that leave the implementer waiting for a framework choice. A Stipula--KeY
experiment remains optional and bounded to one workflow property. If its
abstractions hide the relevant event behavior, report that result and continue
the useful reference implementation.
% END SOURCE UNIT P-04-prototype:02

% BEGIN SOURCE UNIT P-03d-java-delivery:06
\subsection{Release contents, extension tasks and acceptance}

The W-package names below identify the original allocation of responsibilities.
They explain the acceptance obligations that drove the implemented T00--T22
tasks; they are not an additional unfinished backlog. The execution account that
follows maps those obligations to the delivered workbench. The E tasks later in
this chapter specify the new ensemble work.


\label{merge:P-03d-java-delivery:06}


A production change package must contain the approved source and dependency
register, provision dispositions, rule bundle, reviewer records, supported
profile and effective interval, generated Java source, runtime library, binary
JAR, dependency/toolchain manifest, data-mapping contract, tests and verification
report. The report names the exact JAR tested. The bank's deployment process
verifies these digests and authorization, performs its security and change checks,
then switches the active release atomically. A signature can authenticate the
release authority; a digest alone only identifies content.

% BEGIN REVISION ADDITION teach-P-03d-java-delivery-06
\begin{figure}[H]
\centering
\TeachingFlow{Reviewed source and rule}{Tested Java package}{Controlled release decision}
\caption{A release contains evidence of meaning, behavior and authority.}\label{fig:teach-P-03d-java-delivery-06}
\end{figure}

% END REVISION ADDITION teach-P-03d-java-delivery-06
The implementation backlog adds a mandatory Java package, W03J, immediately
after the full reference evaluator. It owns the Java domain/API classes, one
emitter per RuleIR operator, result serialization, exact numeric/date behavior,
compiled conformance harness, JAR manifest and host adapter contract. Its pass
condition is equivalence of the complete supported responses on the full
conformance suite, source-derived scenarios and distinguishing mutations,
plus reproducible generation and a separately compiled caller. A backend that
accepts an operator but approximates its meaning fails the package.

W04 must bind release evidence to both the rule bundle and the resulting JAR.
W05 must extend the demonstration's consent replay to the full event contract,
including obligation instances and the adopted calendar policy. W06 must show
the source, interpretation, test case and Java result together. W10 must include
host concurrency, idempotency, stale release and replay tests in its engineering
checks, then separately measure review effort and material errors in the pilot.
This sequencing produces a Java deliverable early while keeping legal
adjudication and production authority explicit.
% END SOURCE UNIT P-03d-java-delivery:06

% BEGIN SOURCE UNIT P-05-evidence:01
\subsection{Distinguishing specification fixtures from executed evidence}
\label{merge:P-05-evidence:01}


The proposal now has a machine-readable companion to the business scenarios.
Under \path{docs/specs/v0.1/fixtures/}, 35 decision cases specify complete bundles,
synthetic snapshots, assessment times and expected result projections. Six
deliberately invalid variants exercise schema, reference, type and identifier
failures. Eight event histories cover withdrawal, regrant, timely performance,
deadline equality, missing observations and late performance. Four release cases
distinguish valid approval from self-review, open issues and stale evidence.

These sets serve different purposes. The 35 cases specify the selected language;
the fifteen business scenarios above identify regulatory interpretation and
operational questions, several of which remain beyond the first executable
slice. An agent must not claim that passing the language cases completes those
business scenarios. For example, a Boolean assessment input can represent a
reviewer's judgment about investment objectives, but it does not automate that
judgment. Likewise, ownership and category definitions require their own reviewed
normalization and end-to-end tests before release.

% BEGIN REVISION ADDITION teach-P-05-evidence-01
\begin{figure}[H]
\centering
\TeachingCompare{Fixture states expected behavior}{Executed check observes behavior}{Legal review justifies the expectation}
\caption{A fixture file alone is not a passing test or an adjudicated interpretation.}\label{fig:teach-P-05-evidence-01}
\end{figure}

% END REVISION ADDITION teach-P-05-evidence-01
The original specification checker validates three initial JSON schemas,
the 35 decision fixture contracts, six negative examples, event-record structure,
source/derivative/span hashes and the SQLite definitions. Its static mode alone
does not execute the engines. The accepted MVP subsequently executed all 35
decision cases in runtime mode and exercised event and release behavior through
separate tests and the complete walkthrough. The recorded acceptance also checked
regeneration of 31 JSON contracts, including OpenAPI. Static fixture checks and
those later execution results remain different evidence. The separate Java
demonstration now supplies narrower execution evidence: 32 SPI decisions and eleven
consent histories agree with the independently written reference; all three
compiled mutations are detected; a fresh rebuild produces an identical JAR;
and a separately compiled host calls that JAR before and after withdrawal.
The exact commands, toolchain and source/binary hashes are retained in its
verification report. Its ten rules and reviewed-input boundaries are explained
in Sections~\ref{sec:dry-run} and~\ref{sec:java}. Independent legal adjudication,
human usability assessment and published-tool proof replay remain prospective.
% END SOURCE UNIT P-05-evidence:01

% BEGIN SOURCE UNIT P-03c-storage-review:02
\section{The implemented authoring service and its current boundary}

The original service design below names the core responsibilities and transport
conventions. T00--T22 implemented a larger 22-path API; the retained generated
OpenAPI and the exact runtime closure resolve the original sketches. Reviewer
identity is derived from authentication, every mutation uses a caller-scoped
idempotency key, and the build--verification--release sequence is acyclic.
The table is a design summary, not a substitute for the current request models.


\label{merge:P-03c-storage-review:02}


The API imports sources, creates and retrieves bundles, records reviews, releases
approved hashes, evaluates snapshots, appends events, replays histories and
compares versions. The transport contract is summarized below; companion files
preserve the same field names for implementation. The central reference call is

\begin{minipage}{\textwidth}
\begin{lstlisting}
evaluate(bundle, snapshot, rule_id, valid_at, known_at,
         mode="draft")
    -> EvaluationResult
\end{lstlisting}
\end{minipage}

The evaluation response includes the tagged outcome, exact value if any, reason codes,
missing and blocking inputs, trace, source/bundle/input hashes and engine version.
Production calls require a released exact hash; draft calls visibly retain draft
mode. An HTTP failure is different from a well-formed request whose assessment
returns an error or conflict.

The original authoring-service sketch assigns the following endpoint responsibilities. Path parameters name
immutable content identities; state-changing requests also carry a caller-scoped
idempotency key. Repeating the same key and request returns the stored response;
reusing it with different content returns HTTP 409. Optimistic revisions prevent
two reviewers from overwriting each other's changes.
\begingroup\small
\begin{longtable}{P{0.36\textwidth}P{0.54\textwidth}}
\caption{Authoring-service contract. This service is separate from the Java
transaction library.}\\
\toprule Method and path & Required request and returned result \\\midrule
\endfirsthead
\toprule Method and path & Required request and returned result \\\midrule
\endhead
\texttt{POST /v1/sources/import} & Official URL or local blob, reference ID and media type; returns source revision and derivative hashes. Invalid content is 422; fetch failure is 502.\\
\texttt{POST /v1/bundles} & Complete typed bundle; returns immutable hash, draft status and revision 1. Schema or type failure is 422.\\
\texttt{GET /v1/bundles/\{hash\}} & Returns exact bundle, review state and issues; unknown hash is 404.\\
\texttt{POST /v1/bundles/\{hash\}/review} & Decision, comment and expected revision; reviewer identity/role derived from authentication; returns new state/revision. Stale update is 409; invalid transition is 422.\\
\texttt{POST /v1/bundles/\{hash\}/release} & Java-release-manifest hash, expected revision and activation time; returns release record if all guards pass. Missing approval or mismatched code evidence is 422.\\
\texttt{POST /v1/evaluations} & Bundle hash, snapshot, rule ID, mode and both times; returns the complete evaluation. A nonreleased production bundle is 409.\\
\texttt{GET /v1/evaluations/\{hash\}} & Returns preserved request, result and source links; unknown hash is 404.\\
\texttt{POST /v1/event-streams/\{id\}/events} & Event, expected revision and ordering authority; returns stored event/revision. An ID payload mismatch or sequence collision is 409.\\
\texttt{POST /v1/event-streams/\{id\}/replay} & Both times and profile version; returns immutable consent/obligation state with evidence.\\
\texttt{POST /v1/comparisons} & Old/new bundle hashes, rule ID and declared domain; returns a persisted job ID and later a counterexample or a specifically bounded analysis result.\\
\texttt{GET /v1/changes/\{old\}/\{new\}} & Returns changed definitions, rules, sources, affected dependants and evidence that must be regenerated.\\
\bottomrule
\end{longtable}
\endgroup

% BEGIN REVISION ADDITION teach-P-03c-storage-review-02
\begin{figure}[H]
\centering
\TeachingFlow{Authenticated request}{Validated state transition}{Durable response and evidence}
\caption{The authoring API is separate from transaction-time Java execution.}\label{fig:teach-P-03c-storage-review-02}
\end{figure}

% END REVISION ADDITION teach-P-03c-storage-review-02
The persistent relationships follow the questions the reviewer asks. A source
revision points to immutable raw and extracted blobs; an interpretation points
to its spans and unresolved issues; a bundle contains the interpreted rules;
reviews and release records point to that bundle hash. A snapshot identifies
its subject and facts. An evaluation points to the bundle, snapshot and both
times. An event belongs to a stream and has an immutable identifier and ordering
evidence; replay results are separate records. Jobs, idempotency records and audit
events preserve asynchronous work and accepted state transitions. The supplied
SQL represents these relationships with foreign keys and JSON payloads; a Java
release manifest can be retained as an immutable blob referenced from the
release evidence. No mutable row is allowed to change the contents identified
by an old digest.

For content identity, sort JSON object keys by ASCII order, preserve array order
and string code points, use compact UTF-8 and hash the bytes with SHA-256.
Reject duplicate keys, floating-point values, nonfinite values and unpaired
Unicode surrogates. JSON integer metadata is limited to the exactly portable
range $[-(2^{53}-1),2^{53}-1]$; domain integers and money use decimal strings.
This is an explicitly specified local format, not an assertion of conformance
to a different canonical-JSON standard. A result hash excludes itself and
incidental timing/request identifiers. These rules let the Python authoring
service and Java library identify the same inputs and semantic output.

Longer solver or drafting calls use persisted jobs with explicit failed,
cancelled and timed-out states. They cannot alter approval records. Official
documents and model output remain untrusted input: fetching restricts hosts and
redirects, validates content and limits resources; parsing never executes a
document's instructions. The initial service binds to localhost and uses public
sources and synthetic facts. This makes the proposed behavior implementable
before private data access or production integration is considered.
% END SOURCE UNIT P-03c-storage-review:02

\section{What the 22 September local execution established}
\label{sec:executed-workbench}

The original work packages have now produced an engineering MVP. The retained
22 September 2026 acceptance contains 154 passing tests, execution of all
35 RuleIR decision cases, and a complete 23EC35 review-to-Java walkthrough.
Its authority is \texttt{LOCAL\_SYNTHETIC}. This means that named synthetic
identities exercise the review controls on public sources; it does not mean
that the bank has accepted the legal interpretations or deployed the code.
The evidence below is the earlier recorded execution, not a new test run
performed while unifying this manuscript.

Three evidence sets must remain distinguishable. SPI-Demo1 supplies 32 decision
examples and eleven consent histories for its narrower interface. The later
workbench implements the full specified decision language, attributed consent
and the finite achievement-duty profile, with review and release services.
The second-circular exercise then adds 22 named 23EC46 cases, 729 abstract
inputs and four compiled mutants while retaining its DRAFT status. Combining
the documents must not turn those separate denominators into a count of
independently adjudicated legal questions.

The workbench walkthrough imports 23EC35 and both annexes and checks all
51 provision dispositions. The selected business profile is an individual
client in a solicited transaction using execution monitoring under Annex 1
paragraph 8.3(a). Five circulars remain in the retained pilot corpus, but the
implemented SPI profile does not thereby cover every provision in those
circulars. Corporate clients, designated-account monitoring and the other
explicit exclusions retain their own dispositions.

For the original candidate, 32 SPI example outcomes and eight attributed
consent/duty histories run through the exact generated candidate JAR.
Separate synthetic author, meaning-reviewer and engineering-reviewer identities
then approve the exact manifest. A separately compiled Java host invokes the
exported library. A recorded withdrawal changes the selected streamlining
result from true to false. Thus the evidence crosses the actual library
boundary and includes an event that changes a later decision.

A synthetic amendment then raises a financial threshold by one cent. It is an
engineering change example, not an SFC amendment. It receives a new bundle,
build, verification, coverage and scope review, approvals and release. Its
October test interval follows the original September interval. Exporting and
restoring the public/synthetic history reproduces the original September
decision exactly, including its original evidence and result identity. A
successful current evaluation alone would not have established that property.

% BEGIN REVISION ADDITION gap-execution-original
\begin{figure}[H]
\centering
\TeachingFlow{Original September decision}{Synthetic one-cent amendment}{Historical September replay}
\caption{The retained execution checks change and historical reconstruction.}\label{fig:gap-execution-original}
\end{figure}


% END REVISION ADDITION gap-execution-original
The released original class is
\texttt{hk.legalmath.Policy\_87e4416d72408979a341}. Its full content identities
are retained here so that the two releases can be distinguished without a
mutable filename:
\begin{lstlisting}[caption={SHA-256 identities of the retained local walkthrough.}]
Original bundle:
87e4416d72408979a341bdb80581201bbd7306be4bb1d81157693f641025ec4e
Original JAR:
6b069b605a331480d618561b61e26c133f08f657813c779bc8c374a2aea62e51
Synthetic amended JAR:
fb0d1c41a98e8f072834998fbd1c70447a9b54e15198c1d46f683524b5166df2
\end{lstlisting}
The different binary identities are evidence that the release changed; their
meaning comes from the corresponding bundle and verification records.

\subsection{From the original packages to completed engineering tasks}

The W packages describe the design responsibilities. The more detailed T tasks
record their implementation and acceptance. Table~\ref{tab:executed-tasks}
retains every core task and the two distinct outstanding tasks. Its evidence
statements refer to the retained full-suite run; they do not assert that every
listed command was separately rerun.

\begingroup\small
\begin{longtable}{@{}P{0.09\textwidth}P{0.86\textwidth}@{}}
\caption{The completed local workbench and its remaining research tasks.}
\label{tab:executed-tasks}\\
\toprule Task & Behavior established by the retained engineering evidence \\\midrule
\endfirsthead
\toprule Task & Behavior established by the retained engineering evidence \\\midrule
\endhead
T00 & Closed boundary, temporal, identity, event-attribution and release-sequence gaps; fixed schemas and adversarial expectations A01--A16.\\
T01 & Isolated Python 3.11 package, pinned dependencies, CLI and a clean non-editable installation.\\
T02 & Canonical JSON, atomic content-addressed blobs, WAL database, migrations, foreign keys, transactional idempotency/audit and corruption checks.\\
T03 & Source and annex import, raw/text/page/offset verification, extractor identity and original-document display.\\
T04 & Provision dispositions, definition dependencies, applicability dimensions and authenticated meaning review; missing coverage or dependencies blocks release.\\
T05 & Full AST validation, fixed diagnostics, duplicate/reference/type/cycle rejection and bounded expanded call depth.\\
T06 & Dated records, interval-local supersession, late corrections, evidence merging and explicit unknown/conflicting evidence.\\
T07 & All specified RuleIR operators, scope/default/unknown/conflict semantics and all 35 original decision cases.\\
T08 & Independent result/trace checks, including altered links, hashes, order and duplicate rejection, with retained boundary expectations.\\
T09 & Complete specified Java decision semantics, verification through the actual generated class, 32 migrated SPI cases, non-Boolean results and three compiled mutants.\\
T10 & Clean deterministic Java 17 builds, exact-binary evidence and a separate caller requiring no Python or model service.\\
% BEGIN REVISION ADDITION gap-execution-tasks
\bottomrule
\end{longtable}
\begingroup\normalsize
\begin{figure}[H]
\centering
\TeachingFlow{Validate and evaluate typed rules}{Build and check Java}{Review, release and retain history}
\caption{The completed tasks connect semantics to a controlled release.}\label{fig:gap-execution-tasks}
\end{figure}

\endgroup
\begin{longtable}{@{}P{0.09\textwidth}P{0.86\textwidth}@{}}

\toprule Task & Behavior established by the retained engineering evidence \\\midrule
\endhead
% END REVISION ADDITION gap-execution-tasks
T11 & Registry roles, author/reviewer separation, stale-revision rejection, exact-manifest approvals, source/coverage checks and atomic audit.\\
T12 & Legal-interval and operational-time release selection, overlap/retirement rejection, exact-byte export and history restoration.\\
T13 & Attributed event persistence, ordering, completeness, consent generations, duty deadlines, late evidence and explicit leap-day policies.\\
T14 & Complete Python/Java event-result agreement, release verification of the event component and consent-to-decision composition.\\
T15 & Synthetic versioned host with simultaneous reservation, consent/release changes, bounded retries and idempotent orders. Its storage remains in process.\\
T16 & API, local identity, stable errors, persistence, bounded background jobs, cancellation, restart and launch-failure handling.\\
T17 & Source/rule/case/amendment screens and a local API console; Chromium, keyboard, HKD formatting and injection checks.\\
T18 & Bounded comparison with explicit knownness; strict-boundary and OR--AND witnesses replay in both engines; empty-domain, unsupported and timeout outcomes remain distinct.\\
T19 & Restricted discovery, hash-bound pilot manifest, attachment changes, replacement metadata and unsafe-network rejection.\\
T20 & Deterministic drafting provider, exact citations, bounded repair and timeout/manual-review outcomes, with no tool or approval authority.\\
T21 & Definition/source/fact/rule impact, fresh exact-version approvals, historical decisions and unchanged opening versions for existing duties.\\
T22 & Complete original/amended review-to-Java walkthrough, external host, export/restore, full suite, fresh installation and browser checks.\\
T23 & Optional research adapters: not run as part of core acceptance. Real adapter invocation, licensing review and proof-search evaluation remain separate work.\\
T24 & Independent human/model comparative pilot: pending. Requires participants, adjudicated tasks and a registered protocol; no effectiveness result exists.\\
\bottomrule
\end{longtable}
\endgroup

The mapping makes the old backlog usable without creating duplicate work.
W00 corresponds to T01--T02; W01 to T03 and the source-import part of T19;
W02 to T05--T06; W03 to T07--T08; and W03J to T09--T10. W04's coverage,
review, release and amendment responsibilities appear in T04, T11--T12 and
T21. W05 becomes T13--T14; W06 becomes T16--T17; W07 becomes T18.
W08's inventory, discovery and drafting work appears in T04 and T19--T20.
W09 remains the optional T23 work. The engineering part of W10 is exercised
by T15 and T22, while its comparative human/model claim remains T24. This
last split prevents completed software testing from becoming a claim of
measured reduction in compliance errors.

The package layout also became concrete during execution. The Java runtime is
under \path{src/legalmath/java/runtime} so an installed package can generate a
standalone JAR. Export and history restoration are implemented in
\path{src/legalmath/review/releases.py} and
\path{src/legalmath/storage/archive.py}. The synthetic transaction host is
packaged with the Java runtime. These path choices replace the earlier proposed
tree while retaining its acceptance obligations.

\subsection{Reproducing the whole selected change}

The implementation can be exercised without a live drafting model. From a
prepared checkout with the locked environment and retained JDK, the following
command creates a fresh, isolated demonstration. Its output directory must be
empty; it must not overwrite an accepted evidence directory.

\begin{lstlisting}[caption={Reproduce the complete local workbench scenario in a new directory.}]
.venv/bin/legalmath demo --repository . \
  --workdir /tmp/legalmath-unified-demo-new \
  --jdk .localresources/java-toolchain/jdk-17.0.20.1+1
\end{lstlisting}

This command imports the retained sources, checks the selected provision
dispositions, builds and verifies Java, records synthetic review, exports the
JAR, compiles the external host, replays withdrawal, reviews and releases the
synthetic amendment, and exports/restores history. The acceptance condition is
the complete recorded sequence and exact historical replay. A generated JAR
alone is insufficient. The equivalent named scenario is
\texttt{legalmath acceptance --scenario spi-23ec35 --offline --out NEW\_DIRECTORY}.
The listing documents an available command; it was not rerun for this merge.

For a fresh environment, create a Python 3.11 virtual environment. Install the
locked dependencies from \texttt{requirements-dev.lock}, then install the local
package without resolving new dependencies. Use the retained JDK 17 or separately verify a compatible
JDK 17. The package does not require the adjacent research environments. Its
local service exposes five review/operation views and 22 API paths, with a
packaged API console and no CDN dependency. The selected generated Java API is
the String/Map interface explained in Chapter~\ref{ch:verification}.

% BEGIN REVISION ADDITION gap-execution-reproduce
\begin{figure}[H]
\centering
\TeachingFlow{Fresh isolated demonstration}{Inspect exact execution evidence}{Keep prior accepted records intact}
\caption{Reproduction creates new evidence without overwriting the original run.}\label{fig:gap-execution-reproduce}
\end{figure}


% END REVISION ADDITION gap-execution-reproduce
The retained commands for engineering acceptance are:
\begin{lstlisting}
.venv/bin/python -m pytest -q
.venv/bin/python scripts/check_contracts.py
.venv/bin/python scripts/check_spec_pack.py \
  --runtime legalmath.conformance:evaluate_case
python3 scripts/check_execution.py
\end{lstlisting}
The application result record binds the actual prior commands and environment.
API tests used a trusted host context because the agent sandbox blocked a
minimal thread-wakeup probe. Changing pinned dependencies to conceal that
environment difference would not have answered the test question. No full
application suite has been rerun simply because the manuscript was merged.

\subsection{The current service sequence, including release preparation}

The complete service sequence clarifies how the earlier endpoint sketch became
an executable control. Import the source and create a draft bundle. Resolve
its coverage, scope and issues under authenticated identities. Build the
candidate with its decision and event cases. Only then prepare an external
release manifest from the build hash, verification hash, applicability hash
and valid interval. Meaning and engineering approvals bind that exact manifest;
release activation consumes its hash, the expected revision and the activation
instant. This order avoids requiring a report to contain its own future approval.

The implemented paths are shown below. All paths have the prefix
\texttt{/v1}; brace names identify path parameters. Reads require the local
authenticated identity, and mutations additionally require the caller-scoped
idempotency key. Domain records have their own typed validation. Request bodies
cannot grant the caller a role, and unknown fields are rejected.

\begingroup\small
\begin{longtable}{@{}P{0.45\textwidth}P{0.50\textwidth}@{}}
\caption{The current 22-path service, before the proposed ensemble extension.}
\label{tab:current-service}\\
\toprule Method and path after \texttt{/v1} & Operation and decisive input \\\midrule
\endfirsthead
\toprule Method and path after \texttt{/v1} & Operation and decisive input \\\midrule
\endhead
POST \path|/sources/import| & Source ID, official URL, retrieval time, media type and optional uploaded bytes; uploads remain unverified.\\
POST \path|/bundles| & Validate and store the complete immutable rule bundle.\\
GET \path|/bundles/{bh}| & Retrieve the exact bundle and its recorded state.\\
POST \path|/bundles/{bh}/coverage| & Record the provision inventory with an expected revision and authenticated meaning review.\\
POST \path|/bundles/{bh}/issues/{issue_id}| & Resolve an issue with reason, time and expected revision.\\
POST \path|/applicability| & Record the supported business scope and rule-version context.\\
POST \path|/bundles/{bh}/build-java| & Supply decision cases and optional event cases; build and verify the candidate.\\
POST \path|/release-manifests| & Bind build, verification and applicability hashes to a valid interval.\\
POST \path|/bundles/{bh}/review| & Submit, reject, approve meaning or approve engineering; require expected revision and the relevant manifest identity.\\
POST \path|/bundles/{bh}/release| & Activate the exact release-manifest hash at an explicit time under the required approvals.\\
% BEGIN REVISION ADDITION gap-execution-api
\bottomrule
\end{longtable}
\begingroup\normalsize
\begin{figure}[H]
\centering
\TeachingFlow{Prepare the tested release}{Obtain exact-manifest approvals}{Activate the approved version}
\caption{The historical API separates preparation, approval and activation.}\label{fig:gap-execution-api}
\end{figure}

\endgroup
\begin{longtable}{@{}P{0.45\textwidth}P{0.50\textwidth}@{}}

\toprule Method and path after \texttt{/v1} & Operation and decisive input \\\midrule
\endhead
% END REVISION ADDITION gap-execution-api
POST \path|/fact-records| & Retain dated, evidenced facts and explicit supersession records.\\
POST \path|/snapshots| & Resolve a bundle, subject, assessment time and knowledge cutoff into a snapshot.\\
POST \path|/evaluations| & Evaluate the exact bundle/snapshot/rule and both times; production also resolves the release and operational time.\\
GET \path|/evaluations/{ident}| & Retrieve the preserved request and result.\\
POST \path|/event-streams| & Create the attributed stream header before appending events.\\
POST \path|/event-streams/{stream_id}/events| & Append an event under an expected revision; enforce header attribution and ordering.\\
POST \path|/event-streams/{stream_id}/replay| & Replay using both times, profile version and the supplied completeness evidence.\\
POST \path|/comparisons| & Compare old/new bundles over a declared domain and both times, within the bounded solver budget.\\
POST \path|/drafts| & Run the current deterministic candidate fixture with source inventory, up to three responses and at most three attempts.\\
GET \path|/jobs/{ident}| & Retrieve persisted asynchronous status and result references.\\
POST \path|/jobs/{ident}/cancel| & Cancel work without manufacturing a partial successful result.\\
GET \path|/changes/{old_hash}/{new_hash}| & Inspect structural and definition impact; semantic comparison is a separate operation.\\
\bottomrule
\end{longtable}
\endgroup

For example, the current evaluation body names \texttt{bundle\_hash},
\texttt{snapshot}, \texttt{rule\_id}, \texttt{valid\_at},
\texttt{known\_at} and a mode of draft, production or replay. The optional
\texttt{release\_hash} and \texttt{operational\_at} supply release context;
merely sending the word ``production'' does not authorize a bundle. A malformed
canonical input is a boundary failure. A well-formed but semantically invalid
snapshot instead produces the defined error result. This distinction is part
of the interface's observable behavior.

These paths supply the earlier 22-path execution route. The later E01--E09
increment adds interpretation operations with separate records and budgets,
bringing the generated API to 33 paths. The original drafting endpoint
retains its separate deterministic-fixture behavior.

\subsection{What the acceptance checks add, and what they leave open}

Contract regeneration reproduced 31 JSON files, including the current OpenAPI
contract, byte for byte. Separate clean build directories produced identical
JAR bytes, stale injected class files were excluded, and a non-Java-17 toolchain
was rejected. Python/Java comparisons exclude their intentionally different
backend identities while independently reconstructing each native result hash.
The generated policy in the exact JAR is invoked, so comparison does not merely
exercise a generic runner while leaving the generated binding unchecked.

The fresh-install check used a non-editable installation in a new environment.
The browser record covers five Chromium views, an authenticated request, mobile
layout and keyboard expansion of a rule tree. The first browser harness used a
string-evaluation wait that conflicted with the application's content-security
policy. Replacing that wait with locator assertions repaired the harness while
retaining the policy. These checks establish local engineering behavior; they
do not constitute an independent compliance-reader usability study.

The restricted discovery client also recorded 159 product-circular index
entries over two pages in the 22 September refresh. Page identities and cursor
evidence are retained. Tests cover repeated pages, changing totals and fetch
errors. The actual 23EC53/26EC22 replacement relationship is represented, and
25EC48 remains a review announcement without an invented numerical rule. This
is a selected category observation, not the bank's entire regulatory perimeter.

The final engineering repairs addressed validation precedence, hashable malformed
snapshots, exact 6,000-digit values, deep inter-rule references, reviewed provision
coverage, reviewer attribution of initial snapshots, uploaded-source authority
and worker launch failure. These repairs explain why source identity, validation
and failure outcomes are part of observable semantics rather than incidental
implementation details.

The evidence is retained under \path{artifacts/runs/acceptance} and
\path{artifacts/runs/mvp-accepted}. The acceptance manifest binds 119 input
identities; its supplemental record binds 36 evidence files and the documented
browser-harness revision. The walkthrough, build, verification, release and
archive records provide the route from a source packet to the external Java
call and restored historical decision. They remain unmodified by this merge.

% BEGIN REVISION ADDITION gap-execution-limits
\begin{figure}[H]
\centering
\TeachingCompare{Synthetic local identity}{In-memory host integration}{Scripted drafting provider}
\caption{The executed workbench still needs production identity, host and model evidence.}\label{fig:gap-execution-limits}
\end{figure}


% END REVISION ADDITION gap-execution-limits
The operational limits are concrete. Review records and history archives use
unsigned local identities and trust database administration. The transaction
host is an in-memory integration example rather than durable bank order
storage. PDF intake has local file/page/text bounds but no operating-system
isolation for hostile documents. The drafting provider is a deterministic
fixture. No universal compiler proof or new-circular automatic interpretation
accuracy has been established. Those limits determine what the ensemble,
independent evaluation and bank integration still have to supply.

The E01--E14 tasks that follow start from this implemented base. They add an
interpretation investigation service; they do not rebuild the accepted evaluator
or replace the original oracle with generated answers. This makes the full
execution route explicit while leaving legal meaning, empirical effectiveness
and real bank authority as separate acceptance questions.


% BEGIN REVISION ADDITION execution-update
\section{The scripted interpretation increment completed on 23 September}
\label{sec:executed-interpretation}

At the 23 September checkpoint, the increment implemented E01--E09. Its retained final acceptance
records 205 regression tests, all 35 RuleIR conformance cases, regeneration of
42 JSON contracts, a 33-path API, Chromium interaction and an offline installed
wheel check. The phase counts overlap and are not additive. These are the
recorded results in
\path{docs/implementation/interpretation-round1/execution-report.md} and its
hash-bound acceptance files, rather than a new execution of the whole suite
during this document revision.

The distinction from a live interpretation service is concrete. The controller
accepts supplied proposals and a supplied repair through its scripted provider.
It retains source units, alternative families, root issues, action reservations,
terminal reports and exact-version review. It does not discover the correct
English interpretation by a newly evaluated model. General argument adjudication,
paid retrieval, live generation of alternative families and calibrated legal
error probabilities remain outside this increment.

\begin{figure}[H]
\centering
\TeachingFlow{Supplied incomplete proposal}{Controller executes supplied repair}{Java and report checks record the result}
\caption{The executed increment tests the investigation machinery with scripted proposals.}
\label{fig:executed-scripted-increment}
\end{figure}

The demonstration makes that limitation observable. A synthetic candidate omits
the product-type promotional route. Named case \texttt{g02} detects the change;
a supplied repair restores the disjunct. The corrected Java build agrees with
the retained 22 named cases and 729 abstract assignments, and a solver witness
is replayed in Python and Java. This supports the repair's encoding behavior
under the adopted interpretation. It supplies no independent legal judgment.

The real retained 23EC46 candidate keeps its original five legal questions and
two missing Code/FAQ dependencies. Its report therefore carries seven unresolved
material issues, and the candidate remains DRAFT and unreleased. The repaired
synthetic investigation also retains its unresolved meaning and structural
review. A separate fully resolved synthetic fixture exercises the positive
release path with test identities. Combining these outcomes into a claim that
the circular was approved would change their meaning.

\begin{figure}[H]
\centering
\TeachingBranch{Encoding repair passes its cases}{Seven source or legal issues remain}{Release stays blocked}
\caption{Repairing an executable omission does not supply missing authority.}
\label{fig:encoding-repair-still-blocked}
\end{figure}

At this historical checkpoint, restricted providers, reference construction,
comparison and search remained subsequent work. Later rounds implemented live
development generation, bounded search, source checks and local Java integration,
as described in Appendix~\ref{app:current-interpretation}. Fair empirical
evaluation and bank identity, data and operational integration remain unfinished.
The earlier API tables retain their role as design explanations; exact interfaces
must be checked against the generated OpenAPI contract and current command help.

% END REVISION ADDITION execution-update

% BEGIN SOURCE UNIT M08:00
\section{Extend the existing product at a defined boundary}
\label{merge:M08:00}


The current repository already has source intake, RuleIR validation and execution,
event replay, Java generation, review records, build and release manifests, and a
local service interface. Its drafting provider is a deterministic fixture. The
new work is an interpretation investigation service between source intake and
approval of a candidate bundle. Replacing the entire runtime would discard useful
engineering evidence and enlarge the task unnecessarily.

The extension accepts a retained source packet, an activity profile and an
explicit investigation policy. It creates inventories, candidates, arguments,
discrepancies and bounded actions. Its output is a report and, when appropriate,
one or more validated draft bundles. The existing review and release services
remain responsible for approval and export, with an additional requirement that
the interpretation report and its unresolved issues be checked.

% BEGIN REVISION ADDITION teach-M08-00
\begin{figure}[H]
\centering
\TeachingFlow{Existing source and RuleIR services}{Interpretation investigation}{Existing review and release controls}
\caption{The extension adds investigation at a defined product boundary.}\label{fig:teach-M08-00}
\end{figure}

% END REVISION ADDITION teach-M08-00
This boundary has two advantages. The interpretation service can evolve its
candidate-generation methods without changing deterministic production execution.
The Java path can be verified independently of any particular model provider.
A candidate that requires an unsupported operator is returned as an engineering
question rather than forcing a runtime change through an unreviewed model output.

The companion contracts define the \texttt{interpretation.v1} extension.
They are stored in \texttt{docs/monograph/contracts}. E01--E09 now implement
the public-source, scripted increment described below. The retained task
descriptions state its obligations; provider integration, independent
evaluation and bank deployment remain later work. The accompanying ordered task file identifies dependencies and acceptance conditions. An
implementation agent should read this chapter with those files and the existing
v0.1 runtime contracts; it should not infer additional authority from illustrative
code in the manuscript.
% END SOURCE UNIT M08:00

% BEGIN SOURCE UNIT M08:01
\section{Separate durable records from generated prose}
\label{merge:M08:01}


The service needs a small set of durable record types. A run binds the source
packet, activity profile, policy and required ensemble roles. A candidate binds a
proposed meaning to a formal representation and its assumptions. An issue records
a discrepancy and its evidence need. An action records one authorized external
operation. A report reconciles the run's evidence and terminal status. Review
decisions and source coverage records connect these objects to human authority
and document completeness.

Generated prose is a field within these records, not their controlling structure.
A model may write an explanation saying an issue is resolved, but the service
changes the issue's resolution status only through a validated transition with
the required evidence. A sentence saying ``approved by compliance'' has no effect
unless an authenticated review record exists for the exact object.

% BEGIN REVISION ADDITION teach-M08-01
\begin{figure}[H]
\centering
\TeachingCompare{Generated explanation}{Validated candidate record}{Authenticated decision}
\caption{Prose describing approval cannot create approval.}\label{fig:teach-M08-01}
\end{figure}

% END REVISION ADDITION teach-M08-01
Identifiers are stable opaque strings within a namespace; content hashes identify
immutable payloads. These roles should not be confused. A run identifier remains
stable while events accumulate, whereas each immutable candidate revision has a
new content identity. A human-readable title can change without becoming the
primary key. References use identifiers and, where integrity requires it, the
expected content hash.

Timestamps use the project's canonical UTC format. Durations are integer
milliseconds or seconds with explicitly named units. Counts are nonnegative
integers. Monetary reservations use exact minor units and a named billing
currency. No field called \texttt{cost} should leave the unit or estimated-versus-
actual distinction implicit. This discipline prevents the controller from
comparing incompatible budgets.

The schemas reject unknown properties. This catches misspelled limits and
unrecognized status fields rather than ignoring them. Structural validation is
only the first step: the service also verifies references, state transitions,
source compatibility, budget arithmetic and role authority. JSON Schema cannot
prove that a source passage supports an interpretation or that two referenced
objects belong to the same run.
% END SOURCE UNIT M08:01

% BEGIN SOURCE UNIT M08:02
\section{The run and policy records}
\label{merge:M08:02}


A run records its immutable configuration hash, source packet hash, profile
identifier, required roles, creation time, deadline and current lifecycle state.
It also references its candidate, issue, action and report collections. Mutable
projections such as counters are derived from durable events or updated
transactionally with them; they are never trusted solely because a client submits
a number.

The lifecycle begins with initialization, proceeds through initial proposals and
resolution rounds, and ends in ready-for-review, blocked-unresolved, cancelled or
failed-integrity. A ready-for-review run is complete enough to present, not
released. A run that is blocked can still contain useful validated candidates.
The API returns them with the block, rather than suppressing all information
because the overall investigation is incomplete.

The policy supplies all material resource and coverage choices. It names required
roles, maximum initial actions, maximum resolution rounds, per-root and global
action caps, no-progress limit, action timeout, run duration, candidate count,
argument-node count and dependency depth. It also declares whether monetary and
token caps are enforced. Configuration validation checks internal consistency,
such as requiring the initial-action cap not to exceed the global cap.

% BEGIN REVISION ADDITION teach-M08-02
\begin{figure}[H]
\centering
\TeachingFlow{Immutable run policy}{Counters and absolute deadline}{Terminal status with remaining uncertainty}
\caption{A run's bounds remain tied to the configuration under which it began.}\label{fig:teach-M08-02}
\end{figure}

% END REVISION ADDITION teach-M08-02
The reference fixture policy is deliberately explicit:
\begin{lstlisting}[caption={Reference limits for deterministic controller fixtures.}]
{
  "schema_version": "interpretation.v1",
  "policy_id": "fixture-policy.v1",
  "use": "DETERMINISTIC_FIXTURE_ONLY",
  "max_initial_actions": 4,
  "max_resolution_rounds": 3,
  "max_actions_per_issue_root": 3,
  "max_actions_total": 12,
  "max_no_progress_rounds": 2,
  "action_timeout_seconds": 30,
  "run_deadline_seconds": 600,
  "max_candidates": 24,
  "max_argument_nodes": 200,
  "max_dependency_depth": 6,
  "token_cap": null,
  "cost_cap_minor_units": null
}
\end{lstlisting}

The candidate, graph and dependency bounds make local processing tractable for
fixtures. Reaching one opens or preserves an incompleteness issue. These numbers
are convenience choices for engineering tests, not selected production defaults.
A paid provider configuration must add enforceable cost and token reservations,
and its deployment policy must explain why the chosen limits are sufficient for
the intended workload.

A policy version is immutable once used. Changing a cap during a run requires an
explicit authorized amendment event if that capability is later implemented;
the first prototype should instead require a successor run. This keeps budget
semantics simple and prevents a worker from extending its own authority. The
successor links to the prior report and explains why additional work is justified.
% END SOURCE UNIT M08:02

% BEGIN SOURCE UNIT M08:03
\section{The candidate record carries the interpretive commitment}
\label{merge:M08:03}


A candidate has a parent revision, source references, subject unit, controlled-
language statement, formal bundle reference, assumptions and interpretation
family. The family identifies a consequential choice, such as component versus
whole-offer assessment. It is not merely a cluster label assigned by a text
embedding. A candidate can belong to several dimensions, but its changed choices
must be explicit relative to its parent.

Each assumption has a statement, provenance and status. Possible statuses include
supported by inspected source, factual evidence required, provisional modeling
choice, disputed and rejected. A known fact in a runtime snapshot must not be
created directly from a disputed assumption. The candidate-to-fact mapping is a
reviewed transformation with its own evidence requirements.

Supporting and opposing arguments reference source spans or other arguments.
An argument distinguishes premises, inference rule and conclusion, and labels
strict versus defeasible inference where that distinction is used. If an
argumentation procedure is unavailable for a particular structure, these records
still support a structured review interface. Later work implements typed
challenges and bounded argument analysis, without a general legal-priority
procedure. The service must not label the resulting graph as
formally evaluated until the selected semantics has actually been run.

% BEGIN REVISION ADDITION teach-M08-03
\begin{figure}[H]
\centering
\TeachingBranch{Two candidates share one formula}{Different definition assumption}{Different source commitment}
\caption{Formal deduplication must preserve interpretive differences.}\label{fig:teach-M08-03}
\end{figure}

% END REVISION ADDITION teach-M08-03
Candidate status distinguishes active, deferred, defeated, unsupported and
selected-for-review. Deferred means not yet sufficiently investigated, not false.
Defeated requires a recorded reason under the chosen argument and authority
policy. Unsupported means the current record lacks adequate support. Selected
for review is a presentation decision; only a separate authenticated decision can
approve the candidate's meaning.

The record includes score components with definitions and calibration status.
The initial contract permits a scheduling priority and categorical evidence
states. It explicitly leaves probability of legal correctness null. Later
calibration can add a versioned, narrowly defined score, but cannot reinterpret
old uncalibrated priorities as probabilities after the fact.

Source and assumption identity survive formal deduplication. If two candidates
share an identical RuleIR body, the body can be stored once. Their distinct
assumptions and source mappings remain separate candidate records. This prevents
a compiler optimization from erasing the very disagreement the workbench is
intended to expose.
% END SOURCE UNIT M08:03

% BEGIN SOURCE UNIT M08:04
\section{The issue and action records make automatic work accountable}
\label{merge:M08:04}


An issue names its root, parent if any, run, source units, affected candidates,
semantic dimension and materiality. It records the observed discrepancy and the
question that would resolve it. Processing state and resolution state are separate
fields. A terminal processing state can coexist with unresolved meaning.

Resolution evidence is typed. A source-based resolution references retained spans
and the reviewer or deterministic rule that accepted them. An adjudication
references the authenticated decision. An equivalence resolution references the
comparison report, domain and confirmation that no distinct source commitment
remains. A bare model statement is not a valid resolution-evidence type.

An action names the run, issue root or initial role, round, permitted action type,
input hashes, adapter, timeout, reservation and current state. It also stores the
idempotency key, lease owner, fencing token and terminal result reference. The
action's question and success criterion are immutable after dispatch. Changing
them creates another counted action.

The action result separates transport success from substantive progress. A model
request can return valid JSON yet fail to provide any supported new proposition.
A retrieval can succeed technically while returning the wrong version. A solver
can complete normally with an unknown result. The adapter reports its technical
outcome, and the controller validates whether the result answers the action's
question.

% BEGIN REVISION ADDITION teach-M08-04
\begin{figure}[H]
\centering
\TeachingFlow{Issue names missing evidence}{Action asks a bounded question}{Validated evidence permits closure}
\caption{Transport success alone cannot resolve an issue.}\label{fig:teach-M08-04}
\end{figure}

% END REVISION ADDITION teach-M08-04
External text is untrusted data. It cannot name a new tool and cause execution,
change a budget, alter role permissions or select an arbitrary filesystem path.
Permitted actions come from a server-side registry. The model can propose an
action from that registry, but the controller validates its applicability and
reserves resources before dispatch. This is both a security boundary and a way to
keep the method's stopping proof connected to actual behavior.

Typed evidence records bind their content hash, source context, validation state
and origin action or review. Source spans, factual records, validated repairs,
formal comparisons and human adjudications are distinct evidence kinds. A formal
comparison additionally names its domain and result; an unknown result cannot
close an equivalence issue. The companion checker rejects unvalidated closure
and equivalence claims supported only by a repair narrative. Establishing legal
sufficiency and authenticated review authority remains a service and review
responsibility beyond JSON validation.

The issue record retains attempted actions even when they fail. It is often useful
to know that an official dependency was unavailable at the assessment cut-off,
that two classification prompts produced no new evidence, or that a formal check
was outside the supported fragment. Those facts explain why the investigation
stopped and what would justify reopening it.
% END SOURCE UNIT M08:04

% BEGIN SOURCE UNIT M08:05
\section{Coverage records connect source units to consequences}
\label{merge:M08:05}


A coverage record maps each inventoried unit to one or more dispositions:
implemented rule, contextual explanation, definition dependency, separate
obligation, reviewed out-of-scope provision, or unresolved question. The record
includes the reason and responsible reviewer where a legal judgment is involved.
A paragraph can have more than one role; it may contain both a definition and an
operative condition.

Coverage validation first checks that the record refers to the same source hash
as the inventory. It then checks that every inventoried unit appears, every
referenced rule or issue exists, and any claim of implemented coverage has a
validated bundle and verification evidence. It cannot decide the substantive
adequacy of a legal disposition from structure alone, so that judgment remains
an explicit review obligation.

% BEGIN REVISION ADDITION teach-M08-05
\begin{figure}[H]
\centering
\TeachingBranch{Previously missed footnote found}{Official bytes unchanged}{Extraction and coverage need revision}
\caption{An extraction correction can invalidate review without a regulatory amendment.}\label{fig:teach-M08-05}
\end{figure}

% END REVISION ADDITION teach-M08-05
Dependencies are edges with types and versions. An incorporated definition is
stronger than a background analogy in the sense that missing it can directly
prevent interpretation of the selected rule. The impact service uses these edges
to identify affected candidates when a source changes. A generic document list
would not tell the service which conclusions depend on which definitions.

The source inventory itself is versioned and reviewed. An improved extractor may
find a previously missed footnote without any change to the official bytes. That
creates a new derivative and inventory revision, not a new official source
publication. Candidates whose coverage was assessed against the older inventory
need reconsideration. Distinguishing source revision from extraction revision
makes that event explainable.

For the second circular, the executable scope remains paragraph 10. The coverage
record must retain all eighteen paragraphs and four footnotes with their
respective dispositions. A machine-readable count of one implemented paragraph
and seventeen other paragraph dispositions is more honest than a banner saying
``circular translated'' without a denominator. The report can still present the
selected slice as a successful engineering example while showing why whole-
document release is not established.
% END SOURCE UNIT M08:05

% BEGIN SOURCE UNIT M08:06
\section{Storage transactions and event history}
\label{merge:M08:06}


The first implementation can extend the existing transactional store. Tables for
runs, candidates, issues, actions, coverage revisions, reports and review decisions
store canonical payloads with indexed identity and relationship fields. Foreign
keys protect references, unique constraints protect idempotency, and transactions
protect coupled state changes. The exact storage engine can remain the existing
local choice for the prototype; bank deployment requires its own availability and
administrative-control assessment.

The most important transaction reserves an action. It reads the current run and
root counters under a lock or equivalent serializable mechanism, checks all
limits, inserts the action and outbox row, increments counters and appends the
audit event. If any step fails, none commits. A later worker cannot dispatch an
action that lacks the committed reservation.

% BEGIN REVISION ADDITION teach-M08-06
\begin{figure}[H]
\centering
\TeachingFlow{Lock and check remaining budget}{Insert action and increment count}{Commit reservation and outbox}
\caption{Coupled accounting changes belong in one transaction.}\label{fig:teach-M08-06}
\end{figure}

% END REVISION ADDITION teach-M08-06
A resolution transaction similarly checks the issue revision, validates the
resolution evidence, updates the issue projection and invalidates dependent
results when necessary. It does not modify immutable candidate or source blobs.
The event history records what changed and why, while the current projection
supports efficient queries. Recovery can rebuild projections or verify them
against the event sequence.

Client idempotency is scoped by authenticated caller and operation. Repeating the
same key with the same canonical payload returns the existing result. Reusing the
key with a different payload returns conflict. The service must not accept a
client-supplied role in the payload as authority. Authentication determines the
caller and its permitted operations.

The local audit log is useful for reproducibility but is not automatically
 tamper-resistant against a database administrator. A bank deployment may require
external append-only storage, signed checkpoints or equivalent controls. Those
choices are part of the operational security design, not properties conferred by
calling a table an audit log. The prototype should preserve enough identity and
ordering information to integrate such controls later.
% END SOURCE UNIT M08:06

% BEGIN SOURCE UNIT M08:07
\section{An API with reviewable state transitions}
\label{merge:M08:07}


The proposed HTTP extension uses asynchronous jobs for investigation and explicit
resources for review. Creating a run requires the source packet, profile,
configuration and idempotency key. The server validates all references and returns
a run identifier and job reference. It does not block the request until every
model has completed.

\begin{longtable}{@{}P{49mm}P{41mm}P{58mm}@{}}
\caption{Proposed interpretation-service operations. Existing v0.1 endpoints remain separate.}\label{tab:interpretation-api}\\
\toprule
Operation & Purpose & Required behavior\\
\midrule\endfirsthead
\toprule Operation & Purpose & Required behavior\\\midrule\endhead
\texttt{POST}\newline\path|/interpretation-runs| & Create a bounded investigation & Validate source/profile/policy; authenticate caller; idempotent creation\\
\texttt{GET}\newline\path|/interpretation-runs/{id}| & Inspect current state & Return counters, phase, terminal reason and linked resources\\
\texttt{POST}\newline\path|/interpretation-runs/{id}/cancel| & Stop new automatic work & Preserve outstanding actions and produce a cancelled report\\
\texttt{GET}\newline\path|/interpretation-runs/{id}/report| & Retrieve a report version & Expose blocked and failed reports as well as review-ready ones\\
\texttt{GET}\newline\path|/interpretation-issues/{id}| & Inspect a discrepancy & Return candidates, evidence need, attempts and status separately\\
\texttt{POST}\newline\path|/interpretation-issues/{id}/decisions| & Record human adjudication & Require reviewer authority, expected revision and exact evidence links\\
\texttt{POST}\newline\path|/interpretation-runs/{id}/successors| & Authorize further work & Require reason, new policy and lineage; never reset the old run\\
\bottomrule
\end{longtable}

% BEGIN REVISION ADDITION teach-M08-07
\begin{figure}[H]
\centering
\TeachingFlow{Create bounded investigation}{Inspect persistent job and issues}{Retrieve complete terminal report}
\caption{The API preserves an inspectable investigation through asynchronous work.}\label{fig:teach-M08-07}
\end{figure}

% END REVISION ADDITION teach-M08-07
These paths preserve the original design sketch. The executed E01--E09
increment has a generated 33-path OpenAPI contract whose exact request
models govern the implemented interface. Some proposed paths were
consolidated or renamed; a caller must use that generated contract. The book's
path table states responsibilities; the machine-readable schema files state the
record shapes.

Malformed requests return boundary errors without creating partial runs. Missing
references and incompatible source versions return semantic validation errors.
A stale expected revision returns conflict so the reviewer can reload the changed
issue. Unauthorized decisions fail without writing a review or audit entry that
suggests acceptance. The service records denied-operation security evidence
through its appropriate logging channel without pretending that the mutation
succeeded.

Pagination must preserve completeness. Collection responses include stable cursors
and total or explicitly unknown counts; reports bind the complete underlying
collection hashes. A user-interface page showing twenty issues cannot imply that
only twenty exist. Exports include all issue records or a manifest identifying
every page, with integrity checks that detect omission.
% END SOURCE UNIT M08:07

% BEGIN SOURCE UNIT M08:08
\section{The review interface should expose the decisive difference}
\label{merge:M08:08}


A reviewer should begin with the source passage and the decision under review,
not a wall of model transcripts. The interface shows the activity and subject
profile, the candidate's controlled-language rule, and the facts needed to apply
it. A neighboring comparison shows the strongest surviving alternative and a
scenario where the outcomes differ.

The discrepancy view explains the type of uncertainty. Missing factual evidence
should lead to a data request; unresolved meaning should lead to a legal question;
unsupported runtime operations should lead to engineering work. Mixing these
under a single confidence badge makes review slower and encourages the wrong
remedy. The interface can use concise status labels, with the full evidence one
step away.

Source references open the retained version and highlight the exact span. A live
link is useful for checking current authority, but it should not replace the
retained version used in the run. If the live document differs, the interface
shows the change and offers a successor investigation. This prevents a reviewer
from unknowingly approving a candidate against a document that changed after
its creation.

% BEGIN REVISION ADDITION teach-M08-08
\begin{figure}[H]
\centering
\TeachingCompare{Proposed reading}{Strongest supported alternative}{Case where their consequences differ}
\caption{A review screen should make the decisive difference visible.}\label{fig:teach-M08-08}
\end{figure}

% END REVISION ADDITION teach-M08-08
Two review moments require different records. An adjudication during the run
binds the issue revision, candidate and source context; it cannot require a hash
of a final report that does not yet exist. The eventual report includes that
adjudication. A subsequent final review binds the immutable report and candidate.
The companion \texttt{issue-decision} and \texttt{review-decision} schemas keep
these directions separate and avoid a circular report-to-decision hash dependency.
Neither record alone authorizes deployment.

A decision form asks for the proposition being accepted, the evidence relied on,
the treatment of the strongest objection, the scope of the decision and any
conditions. It does not ask only for a confidence slider. The authenticated role
and exact candidate/report hashes are supplied by the service. The reviewer
cannot approve an unseen revision through a stale browser tab.

Unresolved reports remain useful work products. They can present a focused
question for the next authorized reviewer, show which offers are affected and
identify a temporary operational policy proposal. The interface should make clear
whether that policy is a bank choice or a conclusion about the regulator's rule.
This distinction is central to keeping cautious operations from becoming false
statements of law.
% END SOURCE UNIT M08:08

% BEGIN SOURCE UNIT M08:09
\section{Provider adapters are deliberately narrow}
\label{merge:M08:09}


A provider adapter accepts a typed task and returns a bounded response plus usage
and execution metadata. It has no direct database credentials for approvals,
release records or budgets. It cannot invoke arbitrary tools selected in model
text. The controller supplies the exact retained sources and permitted retrieval
scope, and the adapter reports which inputs were actually sent.

The initial interface can support three adapter kinds: scripted fixtures,
retrieval tools and model proposals. Scripted fixtures are essential because they
make failures reproducible. Retrieval adapters return bytes and acquisition
metadata for validation. Model adapters return structured proposals and source
references that the controller checks. All adapters use the same action identity
and timeout accounting.

% BEGIN REVISION ADDITION teach-M08-09
\begin{figure}[H]
\centering
\TeachingFlow{Typed permitted task}{Narrow provider adapter}{Validated response and usage record}
\caption{A model adapter has proposal authority, not release authority.}\label{fig:teach-M08-09}
\end{figure}

% END REVISION ADDITION teach-M08-09
The adapter must distinguish refusal, abstention, malformed output, transport
failure, timeout and unknown remote result. Treating them all as an empty answer
would lose the reason a mandatory role is incomplete. It must also report hidden
provider behavior, such as automatic retries or server-side tool execution, when
that behavior affects resource accounting or evidence provenance. An integration
that cannot bound these behaviors is unsuitable for the unattended profile.

Prompt templates are versioned source files. They state the role, source
perimeter, required structured fields and prohibition on treating source text as
instructions to the system. The prompt should ask for assumptions and opposing
reasons, but those requests do not make the response reliable by themselves.
Validation and independent evidence remain necessary.

The service stores provider and model identifiers, relevant sampling settings,
template version and source hashes. A provider alias that changes its underlying
model must be recorded as such. Reproducibility may be limited even with these
fields, so reports distinguish replay of a retained response from regeneration of
a fresh model response. Deterministic execution of an accepted bundle does not
depend on reproducing the model's original generation.
% END SOURCE UNIT M08:09

% BEGIN SOURCE UNIT M08:10
\section{Security protects the meaning of the evidence}
\label{merge:M08:10}


The source packet can contain malicious or misleading instructions, whether
through a compromised document, an uploaded file or quoted text. The interpreter
must treat them as source content, not commands. A paragraph saying to ignore
prior rules cannot change the controller's budget or authentication policy. The
model's task is to analyze that text within the source context; tool permissions
come from the server.

Acquisition should retain the existing restrictions on HTTPS destinations,
redirects, public addresses, response size and time. A user-supplied URL does not
establish official authority. Uploaded documents remain unverified until a
separate process establishes provenance. A source that is useful for comparison
may still be ineligible to support release.

PDF parsing needs an isolation boundary in a bank deployment. The current local
prototype's bounded parser is not an operating-system sandbox for arbitrary
hostile files. A production intake worker should have limited filesystem and
network access, resource limits and a narrow output contract. This protects the
workbench without changing the legal interpretation method.

% BEGIN REVISION ADDITION teach-M08-10
\begin{figure}[H]
\centering
\TeachingBranch{Untrusted source text says 'approve'}{Read it as document content}{Server keeps approval permissions}
\caption{Source instructions cannot grant tool or reviewer authority.}\label{fig:teach-M08-10}
\end{figure}

A malicious attachment may contain a sentence asking the system to approve its own interpretation. That sentence is evidence to inspect, not an authenticated review. The server must reject an approval operation regardless of how persuasively the document requests it.


% END REVISION ADDITION teach-M08-10
Bank factual data and public regulatory sources have different access needs.
The initial ensemble prototype should use public sources and synthetic scenarios.
Sending client or transaction data to a model provider requires the bank's
approved data-handling arrangement. The architecture supports minimizing that
exposure: legal interpretation can often use abstract fact patterns, while
production evaluation remains inside the bank with the generated Java package.

Logs must avoid becoming an uncontrolled copy of sensitive evidence. Store
references and necessary metadata in operational logs, retain full evidence in
access-controlled stores, and make redaction policies explicit. A redacted report
must preserve enough structure to explain its limitations. Redaction cannot
silently remove the only evidence supporting an approval while leaving the
approval summary apparently complete.

Security tests should use concrete attacks: source text that requests a tool
call, candidate output containing executable markup, a URL redirecting to an
internal address, an oversized graph, a forged reviewer role and an action result
that tries to change its run identifier. The expected behavior is rejection or
safe treatment as data, with a recorded reason. These tests protect the authority
boundaries on which the interpretation method relies.
% END SOURCE UNIT M08:10

% BEGIN SOURCE UNIT M08:11
\section{The ordered implementation sequence}
\label{merge:M08:11}


The implementation should proceed in small stages that each produce a reviewable
capability. The first stage freezes contracts and builds deterministic examples.
The next adds persistence and budgets. Only after the controller behaves correctly
with scripted members should it call live model providers. This sequence avoids
confusing stochastic model behavior with defects in state management.

\begin{longtable}{P{13mm}P{49mm}P{85mm}}
\caption{Original extension work packages at the 23 September checkpoint.
Later generation, search and integration progress is recorded in
Appendix~\ref{app:current-interpretation}; these task descriptions are historical.}\label{tab:extension-tasks}\\
\toprule
ID & Deliverable & Acceptance condition\\\midrule\endfirsthead
\toprule ID & Deliverable & Acceptance condition\\\midrule\endhead
E01 & Contract models and reference policy & All valid examples pass; extra fields, incompatible limits and invalid transitions are rejected.\\
E02 & Source inventory and dependency records & Every retained unit has a disposition; missing or mismatched units create issues even under unanimous proposals.\\
E03 & Immutable candidate and assumption store & Revisions retain parents and source hashes; changing a definition invalidates dependent candidates.\\
E04 & Issue identity and lifecycle & Child issues inherit root accounting; terminal unresolved states cannot masquerade as resolutions.\\
E05 & Atomic budget and outbox service & Concurrent reservations never exceed caps; crash recovery preserves issued counts and fencing.\\
E06 & Scripted provider and controller & Omission repair, persistent ambiguity, timeout, no-progress and cancellation traces match the specification.\\
E07 & Structured comparison and witnesses & Different source contexts are detected; supported formal differences replay in Python and Java.\\
E08 & Complete report and review interface & Every issue and mandatory role is represented; dissent and unexplored families remain visible.\\
E09 & Review and release integration & Exact report/candidate/build identity is required; all stale or unresolved release attempts are rejected.\\
E10 & Restricted retrieval and model adapters & No uncounted retries or model-directed tools; malformed and unavailable responses retain distinct outcomes.\\
E11 & Independent reference corpus & Source and scenario judgments are frozen before model outputs; disagreement and adjudication are retained.\\
E12 & Equal-budget evaluation harness & Baselines share source access and cost accounting; uncertainty and abstention metrics are reported together.\\
E13 & Optional argumentation/search extension & Small reference graphs and family-preservation tests pass; incomplete search remains explicit.\\
E14 & Bank integration and shadow operation & Real identity, host mappings, operational dispositions and rollback procedures are approved and exercised.\\
\bottomrule
\end{longtable}

% BEGIN REVISION ADDITION teach-M08-11
\begin{figure}[H]
\centering
\TeachingFlow{E01–E09 scripted increment}{E10–E12 providers and evaluation}{E14 bank integration}
\caption{Implemented engineering and outstanding empirical and deployment work remain distinct.}\label{fig:teach-M08-11}
\end{figure}

% END REVISION ADDITION teach-M08-11
Dependencies matter. E05 depends on the record and issue contracts, while E06
needs E05 and a scripted adapter. E09 depends on complete reports and existing
release identity checks. E10 can be developed behind the adapter interface, but
live evaluation should wait for E11 and E12. E13 is optional until the simpler
scheduler's limitations are measured; it must not delay a clear implementation
of bounded disagreement handling.

Each task should specify files to add or extend, existing interfaces to reuse,
acceptance identifiers and evidence locations. The companion task JSON provides
that structure without changing the existing master plan's completed T00--T22
history. These are new extension tasks, not a retroactive claim that the earlier
MVP failed to implement its agreed scope.

An implementation agent should stop a task for a genuine contract inconsistency,
missing authority or unavailable required infrastructure. It should not stop
merely because a scripted candidate fails: detecting that failure is often the
acceptance criterion. A failed interpretation candidate can trigger the next
planned repair phase while still blocking release. The task record distinguishes
those outcomes.
% END SOURCE UNIT M08:11

% BEGIN SOURCE UNIT M08:12
\section{Acceptance scenarios that exercise the whole cycle}
\label{merge:M08:12}


The first end-to-end fixture uses the explicit type-of-product omission. It
freezes the source packet, emits two candidates, creates a discrepancy, reserves
a repair action, validates the revised formula, reruns affected checks and
produces a report. The report is ready for review only if all other required
coverage and dependencies in that fixture are complete. The real second-circular
packet remains blocked where its dependencies are unresolved.

The second fixture uses the cashback ambiguity with no resolving authority. It
must preserve both supported candidates through the configured rounds and end in
a blocked report. The test checks not only the terminal status but also the
attempt history, retained objections, question for the reviewer and absence of
release authority. This scenario prevents an implementation from treating
exhaustion as permission to choose the highest score.

The third fixture tests unanimous omission. All candidate providers return the
same incomplete rule, while an independent inventory identifies an uncovered
source unit. The controller must open a material issue despite unanimous agreement.
If it does not, the ensemble has implemented voting rather than the proposed
method.

% BEGIN REVISION ADDITION teach-M08-12
\begin{figure}[H]
\centering
\TeachingCompare{Omission repaired}{Ambiguity persists}{Stale approval rejected}
\caption{Whole-cycle scenarios test different legitimate outcomes.}\label{fig:teach-M08-12}
\end{figure}

% END REVISION ADDITION teach-M08-12
The fourth fixture tests failures under concurrency. Two workers attempt the last
available reservation, one crashes after dispatch, and a stale worker later
returns a result. Exactly one reservation may consume the remaining slot; the
crashed action remains counted; the stale result cannot overwrite a newer state.
The test uses a controlled clock and scripted transport so it is reproducible.
Real distributed-system testing later extends this fixture under the deployment
infrastructure.

The fifth fixture tests invalidation. A reviewer approves a candidate and report,
then a source derivative reveals a previously omitted footnote. The affected
coverage and interpretation records become stale. Release with the earlier
approval must fail until the impact is reviewed and the required checks rerun.
An unchanged Boolean body does not bypass the source-coverage question.

The sixth fixture tests report recovery. A run is cancelled or encounters a
malformed provider response. The service must still produce a complete terminal
report from durable records. If rich rendering fails, a minimal machine-readable
report must remain available. The acceptance condition is not a pretty success
screen; it is a reconstructible account of what happened.
% END SOURCE UNIT M08:12

% BEGIN SOURCE UNIT M08:13
\section{The first vertical slice and its completion rule}
\label{merge:M08:13}


A useful first vertical slice implements E01--E09 with scripted members and the
public second-circular packet. It exercises source inventory, competing candidates,
a repair, a persistent ambiguity, a report and the existing Java path. It does
not need an expensive model integration to test whether the controller's logic
is coherent. This is the smallest product-shaped demonstration of the proposed
method.

Completion requires more than green unit tests. The service must generate an
inspectable report from a fresh run, preserve the exact input and output identities,
show all unresolved issues in the interface, and reject a release attempt for the
blocked example. A separate synthetic fully resolved fixture can demonstrate the
positive review path using local test identities. Its approval status must remain
clearly synthetic.

% BEGIN REVISION ADDITION teach-M08-13
\begin{figure}[H]
\centering
\TeachingFlow{Fresh scripted investigation}{Complete inspectable report}{Independent live-model study still needed}
\caption{The vertical slice demonstrates the controller's specified behavior.}\label{fig:teach-M08-13}
\end{figure}

% END REVISION ADDITION teach-M08-13
The implementation note should record the command, environment, source hashes,
policy, result files and any departures from the design. If a departure changes
meaning, budget accounting or approval behavior, it requires a contract revision
and new acceptance cases. Minor internal choices can be made by the agent, but
must not silently alter observable semantics.

Live models then replace scripted proposal adapters without changing the controller
contract. That transition creates a new empirical question: whether the methods
find useful alternatives and omissions on unseen circulars at acceptable cost.
The next chapter supplies the reference process, baselines and uncertainty
analysis needed to answer that question. Engineering completion of the vertical
slice is a prerequisite for that study, not its result.
% END SOURCE UNIT M08:13


% Detailed comparison tables retained for implementers and auditors.
% The main chapters use compact summaries first and explain how to read them.
\chapter{Detailed comparison tables}
\label{app:detailed-tables}

The tables in this appendix preserve the complete comparison records behind the
shorter tables in Chapters~\ref{ch:languages}, \ref{ch:ensemble} and
\ref{ch:evaluation}. They are lookup material. A reader who wants the argument
can remain in the main chapters; an implementer can use these tables to recover
the full row-by-row specification.

\section{Ensemble responsibilities}
\begin{table}[htbp]
\centering\small
\caption{Responsibilities that create different opportunities to detect error.}
\label{tab:ensemble-roles-detail}
\begin{tabularx}{\textwidth}{P{31mm}YY}
\toprule
Role & Question answered & Failure that must remain visible\\
\midrule
Source inventory & What material is present and incorporated? & Missing attachment, unmatched paragraph, inaccessible dependency\\
Normative reader & What obligations, conditions and exceptions are expressed? & Omitted modality, scope or qualification\\
Alternative reader & What supported choice changes a consequence? & No supported alternative found; unexplored family\\
Argument critic & What defeats the premises or inference? & Unanswered objection or unverified authority\\
Formal checker & What follows from this specified rule? & Invalid encoding, unknown solver result, limited domain\\
Reference reviewer & What is expected for this scenario, and why? & Disagreement, incomplete facts, shared authorship\\
Java checker & Does this package preserve the accepted specification? & Mismatch, unsupported operation, stale build\\
Dossier assembler & Are all required results and dispositions present? & Missing member, open issue, broken evidence link\\
\bottomrule
\end{tabularx}
\end{table}

\section{Resolution evidence}
\begin{longtable}{@{}P{27mm}P{56mm}P{65mm}@{}}
\caption{Discrepancy-specific actions and closure evidence.}\label{tab:resolution-matrix-detail}\\
\toprule
Discrepancy & Appropriate next action & What can justify closure\\\midrule\endfirsthead
\toprule Discrepancy & Appropriate next action & What can justify closure\\\midrule\endhead
Missing source unit & Compare original rendering and independent inventory & Retained unit, corrected inventory and reviewed disposition; agreement among summaries is insufficient.\\
Unavailable dependency & Retrieve the named version or seek authorized clarification & Inspected applicable source, or an explicit adjudication of the remaining question; a guessed definition is insufficient.\\
Different source versions & Establish the assessment cut-off and amendment relationship & A documented version selection and reconsideration of affected candidates.\\
Actor or activity scope & Inspect addressee, operative language and profile facts & Reviewed applicability proposition with evidence for the actual bank activity.\\
% BEGIN REVISION ADDITION gap-resolution-table
\bottomrule
\end{longtable}
\begingroup\normalsize
\begin{figure}[H]
\centering
\TeachingFlow{Classify the discrepancy}{Obtain the relevant evidence}{Close only the question answered}
\caption{Resolution requires evidence of the appropriate kind.}\label{fig:gap-resolution-table}
\end{figure}

\endgroup
\begin{longtable}{@{}P{27mm}P{56mm}P{65mm}@{}}

\toprule Discrepancy & Appropriate next action & What can justify closure\\\midrule\endhead
% END REVISION ADDITION gap-resolution-table
Disputed classification & Identify definition and distinguishing factual evidence & An accepted definition and sufficient facts, or an attributable adjudication; a Boolean label alone is insufficient.\\
Exception attachment & Construct component and whole-arrangement examples & Source-grounded structural choice, with contrary reading disposition and regression cases.\\
Necessary versus sufficient condition & Compare implication directions on boundary scenarios & A reviewed direction of dependence and a rule that preserves other required conditions.\\
Timing or transition & Separate publication, commencement, valid time and knowledge time & Reviewed interval and transition rule with endpoint examples.\\
Factual conflict & Obtain correction or apply an approved evidence-priority rule & Attributable supersession or reviewed normalization; newest arrival alone is insufficient.\\
Formal mismatch & Produce and replay a feasible difference witness & Validated repair or domain-qualified equivalence, with source assumptions retained.\\
Unsupported operation & Inspect whether a supported encoding preserves meaning & Reviewed encoding plus conformance evidence, or an approved runtime extension; compiler rejection does not defeat the reading.\\
Java mismatch & Compare full semantic results and trace & Corrected exact package with repeated affected verification and new build identity.\\
Missing ensemble member & Retry under budget or use an authorized replacement & Required role output with documented independence limits; silence is not agreement.\\
Search truncation & Investigate retained frontier or narrow the claimed scope & Completed declared finite search or explicit reviewed restriction; a high best score is insufficient.\\
Broken approval lineage & Reconstruct exact source, report, bundle and build links & Fresh valid approvals for exact identities; a matching display title is insufficient.\\
\bottomrule
\end{longtable}


\section{Software and standards considered}
\begingroup\small
\begin{longtable}{P{0.19\textwidth}P{0.37\textwidth}P{0.34\textwidth}}
\caption{Candidate software and standards, with a deliberately bounded role.}
\label{tab:software-detail}\\
\toprule Candidate & Contribution to LegalMath & Limit or prototype decision \\\midrule
\endfirsthead
\toprule Candidate & Contribution to LegalMath & Limit or prototype decision \\\midrule
\endhead
Catala & Literate definitions, exceptions, executable decisions and verification-condition ideas. & Compare one adapter on the SPI decision fragment; do not assume event-workflow coverage or end-to-end verified compilation. \citep{catalacode}\\
DMN 1.5; Apache KIE / Drools & Standard decision tables and Java-oriented decision execution. & Preserve hit policies, explicit unknown handling and explanations. Add a separate event model. \citep{dmn2024,droolscode}\\
LegalRuleML & Source, authority, interpretation and normative metadata for interchange. & Adopt a small profile; a reasoner and execution semantics are separate components. \citep{legalruleml2021}\\
Stipula; Stipula--KeY & Contract states, permissions, scheduled events and a Java/JML proof experiment. & Research adapter for a small workflow; document event abstraction and loop restrictions. \citep{stipulacode,stipulakeycode}\\
eFLINT & Normative simulation, duties and violations, including actions that occur despite a prohibition. & Borrow concepts and compare a bounded scenario; pin the implementation and semantics version. \citep{eflint2026}\\
L4 & Legal DSL, regulative rules, event traces, IDE and generated decision services. & Candidate for a later authoring/core comparison; documentation does not establish equivalence with RuleIR. \citep{l4code}\\
Symboleo & Obligation/power lifecycles and a model-checking toolchain. & State-model alternative; preliminary paper case and software inspected, later journal paper remains a retrieval gap. \citep{symboleo2020,symboleocode}\\
Regorous / PCL & Defeasible norms connected to annotated business processes. & Pilot evidence supports the method; deployed version, availability and license remain adoption questions. \citep{semanticcompliance2016}\\
Blawx / s(CASP) & Visual rule authoring, explanations and hypothetical reasoning. & Authoring experiment; Blawx describes its current scope as educational and experimental. \citep{blawxcode,slaw2024}\\
Z3; existing Lean route & Counterexamples, satisfiability, selected invariants and checked proof targets. & Specify supported theories and assumptions; preserve timeout and unknown outcomes. \citep{z3code,localrepos}\\
OpenFisca & Versioned computable rules and a tax/benefit microsimulation architecture. & Useful design analogue; the primary problem here includes conduct and workflow duties. \citep{openfiscacode}\\
Open Policy Agent & Rego policies evaluated against structured data; potential deployment adapter. & Its undefined-value behavior must be reconciled with the legal rule semantics. It supplies no natural-language legal interpretation. \citep{opadocs}\\
Accord Project / Cicero & Adjacent work on digital contracts and reusable contract technology. & Monitor for document/template integration; a fixed contract template differs from an amended regulatory corpus. \citep{cicerocode}\\
\bottomrule
\end{longtable}
\endgroup


\section{Cross-circular acceptance cases}
\begingroup\small
\begin{longtable}{P{0.09\textwidth}P{0.41\textwidth}P{0.4\textwidth}}
\caption{Initial acceptance scenarios for the selected scope.}
\label{tab:cases-detail}\\
\toprule Case & Synthetic fact or source change & Required demonstration \\\midrule
\endfirsthead
\toprule Case & Synthetic fact or source change & Required demonstration \\\midrule
\endhead
F1 & Portfolio HK\$40m; qualifying net assets HK\$0. & Establish the financial
condition via the inclusive portfolio comparison.\\
F2 & Portfolio HK\$35m; qualifying net assets HK\$90m. & Establish the financial
condition via the alternative net-asset route.\\
F3 & Portfolio unknown; qualifying net assets HK\$70m. & Leave the financial
condition unresolved; identify missing portfolio evidence.\\
F4 & Portfolio unknown; qualifying net assets HK\$80m. & Establish the financial
condition via the known sufficient route.\\
F5 & Assets HK\$120m, liabilities HK\$15m, primary residence HK\$30m. & Derive
HK\$75m for the illustrated net-asset calculation; assess the separate portfolio
route independently.\\
O1 & HK\$60m joint portfolio with a non-associate; written client allocation 20\%. &
Use HK\$12m for this portfolio-share illustration; explain the agreement and actor roles.\\
K1 & Transaction-experience route established only for one category; a different
category proposed. & Do not propagate a global experience flag; inspect the
alternative qualification routes and selected categories.\\
C1 & A qualifying wealth record, but conservative investment objectives. & Explain
the investment-objective exclusion; wealth alone does not establish SPI treatment.\\
E1 & Designated-account monitoring selected; exposure rises above threshold without
new funds. & Show the relevant FAQ conditions and restrictions; avoid a universal
transaction-block rule.\\
E2 & A withdrawal instruction precedes a proposed streamlined decision in the
adopted event order. & Use the changed consent state and retain the original
instruction and ordering evidence.\\
T1 & A third-party verification clause with an SFC-request trigger. & Preserve the
trigger; distinguish an absent request from missing request records.\\
T2 & Old tokenisation source replaced by the 2026 circular. & Preserve both editions,
record replacement, identify affected rules and reproduce historical decisions.\\
M1 & Product-linked gift versus a discount of fees or charges. & Preserve the stated
exception and the scope of the cited marketing provision.\\
M2 & Daily dealing fund with weekly redemption versus a different daily cut-off. &
Distinguish reduced dealing frequency from administrative procedure.\\
D1 & A thematic-review announcement enters the source collection. & Classify and
track the announcement without inventing a new numerical eligibility criterion.\\
\bottomrule
\end{longtable}
\endgroup


\section{Evaluation evidence contract}
\begingroup\small
\begin{longtable}{P{0.25\textwidth}P{0.39\textwidth}P{0.26\textwidth}}
\caption{Proposed evidence contract for the prototype evaluation. These are
prospective criteria, not results.}\label{tab:evaluation-detail}\\
\toprule Evidence & Definition and use & Decision role \\\midrule
\endfirsthead
\toprule Evidence & Definition and use & Decision role \\\midrule
\endhead
Material interpretation and decision defects & Independently adjudicated errors
in source coverage, rule meaning, resulting decisions or obligations. No unresolved
critical defect may be carried into a proposed operational scope. & Primary quality
criterion; blocks promotion of the affected scope.\\
Review effort & End-to-end person-time for a completed, accepted change, including
source work, corrections, escalations and reusable-rule preparation. & Primary
business criterion, conditional on quality.\\
Unknown and review outcomes & Frequency, reason, downstream handling and time to
resolution; distinguish necessary abstention from unhelpful failure. & Explanatory
diagnostic and business cost; incorrect silent resolution is a quality defect.\\
Source and dependency completeness & Every in-scope rule has inspected support and
resolved imports; every required selected clause is accounted for. & Release veto
for the selected corpus. Does not establish completeness of all HK regulation.\\
Adapter agreement and mutations & Supported semantics agree; targeted changes in
operators, scope, dates and exceptions are detected by the required checks. &
Engineering veto on the adapter; mutation scores alone do not establish legal fidelity.\\
Reproducibility & A saved rule version, fact snapshot and event sequence reproduce
the original result and explanation references. & Engineering veto on historical replay.\\
Latency and token/runtime cost & Measure on the specified local workload, including
failed drafts and review retries. & Explanatory and operating-budget evidence; no
quality trade-off is implied.\\
\bottomrule
\end{longtable}
\endgroup


\section{Operational readiness evidence}
\begin{longtable}{@{}P{31mm}P{60mm}P{57mm}@{}}
\caption{Evidence required for a bounded deployment decision.}\\
\toprule
Area & Evidence to inspect & Question that remains a human responsibility\\\midrule\endfirsthead
\toprule Area & Evidence to inspect & Question that remains a human responsibility\\\midrule\endhead
Meaning & Source inventory, candidates, objections, adjudication & Is this reading justified for the selected activity and time?\\
Data & Fact dictionary, provenance, validity, conflict policy & Do the bank's records establish the facts the rule needs?\\
Execution & Conformance, witnesses, mutations, Java comparison & Is the supported implementation scope adequate for the process?\\
Authority & Exact manifests, reviewer roles, signatures & Are the required people authorized and accountable for this release?\\
Operation & Host dispositions, race tests, fallback and incident rehearsal & Can staff act correctly when the result is unknown or blocked?\\
Evidence & Retained sources, histories, reproducible reports & Can a later reviewer reconstruct the decision without relying on memory?\\
\bottomrule
\end{longtable}

\section{Recording a method-comparison decision}
\label{sec:comparison-decision-record}

The main volume explains how a comparison affects a decision. This table retains
the fields needed to record the result consistently.

\begin{table}[htbp]
\centering\small
\caption{Decision record required after a substantive method comparison.}
\begin{tabularx}{\textwidth}{P{39mm}Y}
\toprule
Decision field & Required interpretation\\\midrule
Primary criterion & State the observed result and uncertainty against the frozen criterion.\\
Veto status & Identify severe errors, integrity failures and invalid comparisons before discussing speed.\\
Viable candidates & List methods that remain eligible for further evaluation; do not imply an unsupported ranking.\\
Main uncertainty & Distinguish reference ambiguity, sampling error, distribution shift and untested implementation.\\
Next justified action & Promote a bounded use, extend the study, repair a mechanism, or rebuild an invalid harness.\\
Limit of conclusion & State the source perimeter, task population and deployment claims not established.\\
\bottomrule
\end{tabularx}
\end{table}

% Reproducibility details; the main argument does not depend on reading this appendix.
\chapter{Software interfaces and reproduction notes}
\label{app:technical-reference}

This appendix records the concrete software interfaces and historical examples
behind the main narrative. It is intended for someone reproducing or extending
the implementation. The mathematical definitions, worked decisions and proof
assumptions remain in the main chapters.

% BEGIN SOURCE UNIT P-03d-java-delivery:00
\section{The earlier SPI compiler: a concrete deployment lesson}

The API and commands in the next five subsections belong to the retained
\emph{SPI-Demo1} teaching example under \path{examples/java-dry-run/}. They remain
reproducible, but are not the current v0.1 SDK. The subsequent section on the actual
generated interface supplies the current String/Map API. Keeping both makes the
translation mechanism inspectable without directing an implementer to the wrong
library.


\label{merge:P-03d-java-delivery:00}

\label{sec:java}

Java is the required deployment output. The authoring workbench may use Python
for source processing, review and reference evaluation; a transaction must be
able to call the released Java library without running Python or contacting a
language model. The initial integration assumption is a plain Java 17 JAR.
This keeps the decision code independent of Spring, application servers and
particular rule engines. A bank on another Java version needs a declared target
and the same conformance checks before receiving a compatible build.
% END SOURCE UNIT P-03d-java-delivery:00

% BEGIN SOURCE UNIT P-03d-java-delivery:01
\subsection{SPI-Demo1: the narrower retained demonstration}
\label{merge:P-03d-java-delivery:01}


The delivered \texttt{SPI-Demo1} example contains a source-anchored RuleIR bundle,
synthetic fact snapshots, a separately written Python evaluator, a deterministic
Java emitter, the generated class, two small Java support classes, and a separate
host caller. Its implemented expression subset is deliberately explicit:
Boolean, integer and HK-cent literals and facts; unconditional rule references;
\texttt{all}, \texttt{any}, \texttt{not}; and exact numeric
\texttt{eq}/\texttt{ge}/\texttt{gt} comparisons. The root has a Boolean scope.
The demo compiler rejects dates, arithmetic expressions, general defaults,
conditionals and the full obligation language. The later RuleIR implementation
supports the specified decision operators and its finite duty profile; those
features are outside this earlier demo API. The demonstration response is its own smaller API profile,
not a claim to implement every field and trace node of the full RuleIR service.

% BEGIN REVISION ADDITION teach-P-03d-java-delivery-01
\begin{figure}[H]
\centering
\TeachingCompare{SPI-Demo1's limited operators}{Full RuleIR runtime}{Different API and evidence sets}
\caption{The early compiler example and the later runtime have distinct scopes.}\label{fig:teach-P-03d-java-delivery-01}
\end{figure}

% END REVISION ADDITION teach-P-03d-java-delivery-01
The distinction matters for an implementing agent. The example demonstrates an
actual compiler and integration path that can be extended. It does not silently
turn the 35 general RuleIR contract fixtures into executed tests. Those fixtures were the completion target for the full interpreter and
backend and were subsequently executed by the accepted MVP. The separate
SPI example has 32 executed decision scenarios and eleven consent histories, with
its own recorded evidence.
% END SOURCE UNIT P-03d-java-delivery:01

% BEGIN SOURCE UNIT P-03d-java-delivery:02
\subsection{The SPI-Demo1 input and output boundary}
\label{merge:P-03d-java-delivery:02}


The public API accepts an immutable \texttt{Snapshot}, a time being assessed and
a knowledge cutoff. A \texttt{Snapshot} binds a subject identifier to the
declared fact map. A \texttt{Fact} carries a declared scalar type and one of
\texttt{known}, \texttt{unknown} or \texttt{conflict}. A known Boolean is a Java
\texttt{Boolean}; a known monetary or count value is a \texttt{BigInteger}.
Its evidence identifiers, validity interval and recording time are mandatory.
Unknown has a reason and no value. Conflict identifies competing evidence.

The host supplies every declared field, using an explicit unknown where evidence
is unavailable. Omitting a field or supplying the wrong type is a malformed
request, not a legal conclusion. Validity intervals are half-open: evidence
valid until 10:00 does not establish a fact at 10:00. Evidence recorded after
the knowledge cutoff is unavailable to that replay. The API takes canonical UTC
timestamps with six fractional digits; conversion from a Hong Kong business
timestamp belongs in the host adapter and must preserve the intended instant.

\begin{lstlisting}[language=Java,caption={The library API and one dated monetary input.}]
DecisionRuntime.Result evaluate(
    DecisionRuntime.Snapshot snapshot,
    String validAt,
    String knownAt);

// A known value is exact and supported by dated evidence.
Fact portfolio = Fact.known(
    "money_hkd", new BigInteger("4000000000"),
    List.of("valuation.lee", "ownership.lee"),
    "2026-09-01T00:00:00.000000Z",
    "2026-10-01T00:00:00.000000Z",
    "2026-09-01T00:00:00.000000Z");

Fact missingPortfolio = Fact.unknown("money_hkd", "MISSING");
\end{lstlisting}
The first amount is HK\$40 million expressed in cents. A bank adapter reading a
decimal amount should specify the currency, use exact decimal arithmetic and
reject an unsupported fraction of a cent after the adopted conversion. In Java,
\texttt{amount.movePointRight(2).toBigIntegerExact()} can enforce exact cents
after that policy has been applied. Calling \texttt{doubleValue()} first would
discard the property the specification is trying to preserve. For a currency
conversion, the exchange-rate source, date, direction and rounding rule are part
of the evidence mapping.

% BEGIN REVISION ADDITION teach-P-03d-java-delivery-02
\begin{figure}[H]
\centering
\TeachingFlow{Exact dated Snapshot}{SPI-Demo1 evaluation}{Tagged result and disposition}
\caption{The earlier Java interface illustrates the input and output boundary.}\label{fig:teach-P-03d-java-delivery-02}
\end{figure}

% END REVISION ADDITION teach-P-03d-java-delivery-02
The result contains a logical status and an operational disposition:
\begin{center}\small
\begin{tabularx}{\textwidth}{P{0.18\textwidth}P{0.34\textwidth}Y}
\toprule Status & Disposition & Host meaning \\\midrule
\texttt{TRUE} & Conditions met & Selected streamlining conditions established; continue the other controls.\\
\texttt{FALSE} & Standard process required & At least one required condition is refuted; do not apply this streamlined route.\\
\texttt{UNKNOWN} & Review required & Evidence can change the conclusion; show blocking inputs.\\
\texttt{CONFLICT} & Review required & Resolve inconsistent evidence; preserve both records.\\
\texttt{OUT\_OF\_SCOPE} & Outside profile & Select another implemented profile or the standard process.\\
\bottomrule
\end{tabularx}
\end{center}
The demo's JSON includes \texttt{profile}, \texttt{authority}, both timestamps,
the bundle and snapshot hashes, a postorder rule trace, missing and blocking
input lists, and the result hash. A rule-trace entry gives its identifier,
status, source-anchor identifiers and blocking inputs. Source anchors resolve
through the bundle included with the release. Known values and their evidence
remain in the preserved snapshot. The current RuleIR implementation provides expression-node traces,
standardized diagnostics and engine identity through the different v0.1 API
documented later in this chapter.

A malformed request throws an explicit Java input exception in this small
library. The production adapter must map that exception, a timeout or an
unavailable release to a technical-failure result and a review/standard-process
path. It must never catch an exception and substitute a successful Boolean.
The demonstration result always carries \texttt{DEMONSTRATION\_ONLY}; production
authority requires a different, approved release record and enforcement by the
host deployment process.
% END SOURCE UNIT P-03d-java-delivery:02

% BEGIN SOURCE UNIT P-03d-java-delivery:03
\subsection{How the compiler emits final Java rather than a suggestion}
\label{merge:P-03d-java-delivery:03}


The emitter first validates the JSON schema, references, identifiers, types and
acyclic rule graph. It then assigns an internal Java method name to each rule
and recursively emits its expression. Identifiers are validated data, never
pieces of executable text supplied by a language model. The generated monetary
comparison is represented by actual Java calls such as:
\begin{lstlisting}[language=Java,caption={The generated financial method's essential body. The complete class is built and tested in the supplied example.}]
return c.rule("spi.financial", sourceIds, () ->
    any(
        compare("ge", c.fact("portfolio"),
            Value.of(new BigInteger("4000000000"))),
        compare("ge", c.fact("net_assets_ex_home"),
            Value.of(new BigInteger("8000000000")))
    ));
\end{lstlisting}
\texttt{compare} returns a known Boolean when both operands are known, or an
unknown with its blocking input names. \texttt{any} implements the truth table
in Equation~\eqref{eq:or}. Java evaluates its arguments before calling the
method, so both branches have the specified eager evaluation and evidence
collection. Replacing this with Java's short-circuit \texttt{||} would change
the trace and missing-input behavior even where the final Boolean agreed.
The context memoizes referenced rules and records their results in a stable
order. A precheck detects conflicting facts in this profile's dependency set
before normal expression evaluation.

% BEGIN REVISION ADDITION teach-P-03d-java-delivery-03
\begin{figure}[H]
\centering
\TeachingFlow{Typed comparison in RuleIR}{BigInteger comparison in Java}{Same inclusive integer ordering}
\caption{The monetary comparison preserves meaning only if the input mapping preserves units.}\label{fig:teach-P-03d-java-delivery-03}
\end{figure}

% END REVISION ADDITION teach-P-03d-java-delivery-03
The production emitter should implement one mapping per supported operator,
with an explicit error for every unsupported construct. It must not ask an LLM
to write new Java on each regulatory change. The LLM can propose the typed
specification; after review, a deterministic compiler performs the mechanical
translation. This concentrates semantic review on a small language and makes
the generation step reproducible.

For the money comparison, the preservation argument is short. The IR represents
both amounts as mathematical integers of HK cents. The Java boundary constructs
the same integers as \texttt{BigInteger} values. The emitted comparison returns
true exactly when \texttt{compareTo} is nonnegative, which expresses the same
inclusive ordering. The argument depends on the input mapping preserving those
integers. It would fail if the adapter truncated decimals, mixed currencies or
overflowed a narrower integer before constructing the value. For a complete
compiler, analogous local arguments must be composed for every operator,
unknown/error case, rule reference and trace requirement. The delivered tests
check examples of preservation; they are not that universal proof.
% END SOURCE UNIT P-03d-java-delivery:03

% BEGIN SOURCE UNIT P-03d-java-delivery:04
\subsection{Reproduce the build and see the host result}
\label{merge:P-03d-java-delivery:04}


The repository retains a project-local Temurin JDK 17 for this check, verified
against the official archive checksum. It does not change the machine's default
Java installation. From the repository root, the single build command is:

\begin{minipage}{\textwidth}
\begin{lstlisting}
python3 examples/java-dry-run/build_and_verify.py \
  --jdk .localresources/java-toolchain/jdk-17.0.20.1+1
\end{lstlisting}
\end{minipage}

The Python authoring check uses the already available JSON Schema package. The
generated Java library itself has no third-party runtime dependencies. The build
uses \texttt{javac --release 17 -Xlint:all -Werror}; creates the JAR; compiles
the separate fixture and host callers against that JAR; compares complete
reference and Java results; and runs the three mutations. The package has fixed
ZIP timestamps to make repeated builds with the same toolchain reproducible.
The verification record stores compiler identity, source/compiler/JAR digests,
the executed command, results and limitations.

% BEGIN REVISION ADDITION teach-P-03d-java-delivery-04
\begin{figure}[H]
\centering
\TeachingFlow{Compile source and runtime}{Verify the exact JAR}{Call it from a separate host}
\caption{The demonstration crosses the actual Java library boundary.}\label{fig:teach-P-03d-java-delivery-04}
\end{figure}

% END REVISION ADDITION teach-P-03d-java-delivery-04
The delivered files below are operational copies of the design, not additional
reading needed to understand why it works:
\begin{center}\small
\begin{tabularx}{\textwidth}{P{0.45\textwidth}Y}
\toprule File under \path{examples/java-dry-run/} & Purpose \\\midrule
\path{spec/spi-control.bundle.json} & Source-linked typed interpretation for the chosen profile.\\
\path{spec/decision-cases.json}; \path{spec/consent-cases.json} & Independent expected statuses and synthetic evidence/history inputs.\\
\path{reference.py}; \path{compile_java.py} & Separate reference evaluation and deterministic Java generation.\\
\path{generated/GeneratedSpi.java} & Final generated decision class.\\
\path{src/main/java/hk/legalmath/spi/} & Evidence, result, three-valued arithmetic/logic support and consent replay.\\
\path{build/spi-controls-demo.jar} & Compiled library callable by a Java host.\\
\path{generated/BankHostExample.java} & Complete synthetic host caller using the public API.\\
\path{build/verification.json}; \path{build/host-output.txt} & Executed validation and host results.\\
\bottomrule
\end{tabularx}
\end{center}

The separate host prints these dispositions for the synthetic history:
\begin{lstlisting}
BEFORE_WITHDRAWAL STREAMLINING_CONDITIONS_MET
AFTER_WITHDRAWAL STANDARD_PROCESS_REQUIRED
RESULT_AUTHORITY DEMONSTRATION_ONLY
\end{lstlisting}
This is the concrete answer to how final code reaches Java: the bank depends on
the released library, constructs the declared evidence snapshot through its
adapter, calls the generated class and handles its structured result. The
compiler, public circulars and drafting tools are build-time/review-time
components. They need not be deployed into the bank's order-processing process.
% END SOURCE UNIT P-03d-java-delivery:04

% BEGIN SOURCE UNIT M07:09
\section{The actual generated Java interface}
\label{merge:M07:09}


The current emitter validates the bundle, canonicalizes it and embeds its exact
bytes in a generated class. The class name is
\texttt{hk.legalmath.Policy\_} followed by the first twenty hexadecimal digits of
the bundle hash. The class exposes the full bundle hash and delegates evaluation
to the Java semantic runtime. This is a generated, immutable policy binding; it
is not the older proposal's imagined record-based API or a proof that each legal
rule has been translated to hand-optimized Java control flow.

The two current entry points are:
\begin{lstlisting}[language=Java,caption={Implemented generated-policy API.}]
public static String evaluate(
    String snapshotJson,
    String ruleId,
    String validAt,
    String knownAt,
    String mode);

public static Map<String,Object> evaluate(
    Map<String,Object> snapshot,
    String ruleId,
    String validAt,
    String knownAt,
    String mode);
\end{lstlisting}

The string overload returns strict JSON. The map overload returns a deeply
immutable result. The generated policy and Java runtime operate without Python,
a language model, a solver or network retrieval. This property is important for
a bank integration: the interpretation workbench prepares and verifies a package,
while production decision execution can remain local and deterministic.

A concrete host compiles against the actual generated class name from the build
manifest. It supplies a validated snapshot, exact rule identifier, both assessment
times and the requested mode. The mode does not grant authority. A caller cannot
turn a draft into an approved control by passing the word ``production.'' The
release service and host must establish approval and applicability separately.

% BEGIN REVISION ADDITION teach-M07-09
\begin{figure}[H]
\centering
\TeachingFlow{Generated class binds bundle bytes}{String or Map entry point}{Immutable semantic result}
\caption{The current generated policy delegates to the verified runtime.}\label{fig:teach-M07-09}
\end{figure}

% END REVISION ADDITION teach-M07-09
The conceptual integration looks as follows. The symbolic class name in the
listing is replaced by the actual generated name when the integration fixture is
built; the release lookup and transaction methods are host responsibilities.
\begin{lstlisting}[language=Java,caption={Host responsibilities around deterministic evaluation.}]
ApprovedRelease release = releases.resolve(profile, validAt, knownAt);
release.verifyAuthorityAndDigests();
ConsistentSnapshot s = adapter.readSnapshotWithVersions(subject, order);
String resultJson = Policy_HASH.evaluate(
    s.strictJson(), release.ruleId(), validAt, knownAt, "production");
Decision decision = decoder.decodeStrictly(resultJson);
HostDisposition disposition = hostPolicy.map(decision);
transactions.commitIfVersionsUnchanged(
    s.versions(), order.id(), order.payloadHash(),
    release.identity(), decision, disposition);
\end{lstlisting}

\texttt{ApprovedRelease}, \texttt{ConsistentSnapshot} and the host transaction
interfaces in this illustration are proposed integration concepts, not classes
already supplied by the generated SDK. Their responsibilities must be implemented
in the bank's environment. The example makes the boundary explicit so an agent
does not mistake illustrative types for an existing dependency.

The host must handle every tagged result. A true selected prohibition condition
may require stopping the relevant offer; false permits only progression to other
controls under the host policy. Unknown, conflict, out-of-scope and error each
have explicit dispositions. No result string should be cast to a Boolean by a
language's truthiness rules. The decoder rejects unknown status tags so that a
future runtime extension cannot silently take an old default branch.
% END SOURCE UNIT M07:09

% BEGIN SOURCE UNIT M05:10
\section{What can be reused from the adjacent projects}
\label{merge:M05:10}


DynareMCP contains useful patterns for explicit hypotheses, diagnostics,
falsifiers and limits of claims. In the inspected legacy audit writer,
\srcpath{write_hypothesis_node} records a hypothesis, source support, code support,
a diagnostic, a falsifier and a supplied rank score. These fields map naturally
to interpretation candidates. The supplied score is not automatically calibrated;
its meaning depends on the caller.

The inspected legacy scheduler chooses among eligible frontier nodes by a stored
priority and identifier. It is not a general UCT engine hidden behind the word
``tree.'' Reusing its bookkeeping ideas is sensible; importing it and claiming
that a proven Monte Carlo search now verifies legal meaning would be wrong. The
new controller needs its own issue identity, source lineage, legal review states
and resource accounting.

% BEGIN REVISION ADDITION teach-M05-10
\begin{figure}[H]
\centering
\TeachingFlow{Adjacent project's function}{Explicit local assumptions}{Bounded adapter check}
\caption{Reuse must name a concrete capability and the changed context.}\label{fig:teach-M05-10}
\end{figure}

% END REVISION ADDITION teach-M05-10
MathDevMCP is useful once a proposition is precise enough to prove. A formal
obligation might concern preservation of a Boolean expression, absence of an
unbounded loop in a finite controller model, or equivalence of two encodings over
a declared domain. ResearchAssistant supports retrieval, parsing and retention
of the literature and regulatory evidence. Neither project supplies a privileged
interpretation of English. They strengthen different links in the chain.

The design that follows therefore combines several structures: a source inventory
to detect missing material, a candidate set to expose choices, an argument graph
to record reasons, a bounded search to investigate questions, and formal checks
to compare executable consequences. The diversity is useful because each
structure can reveal mistakes that the others overlook. Its value will depend on
how discrepancies move between them, which is the central subject of the next
chapter.
% END SOURCE UNIT M05:10

% BEGIN SOURCE UNIT P-04-prototype:03
\subsection{Adapter boundaries and the earlier reuse assessment}

The earlier repository inspection contributes the adapter requirements below.
The preceding analysis identifies the actual search-controller differences;
references here to historical UCB-style investigations do not override that
finding or establish that a generic MCTS service is already available.


\label{merge:P-04-prototype:03}


The adjacent projects supply reusable patterns and potential adapters. They are
not interchangeable components of a legal reasoner. Their local code and
documentation were inspected at the revisions recorded in the source manifest
\citep{localrepos}. The proposal makes no change to those repositories and does
not assume that their current interfaces implement the proposed legal semantics.

ResearchAssistant is useful immediately for source intake, local storage, parsing,
discovery and structured review records. It was used in this investigation to
discover literature and parse retained PDFs. The available extraction route used
\texttt{pdftotext} and reported low confidence requiring manual review. LegalMath
needs a regulatory layer above it: circular identifiers, annex relationships,
paragraph anchors, versions, replacement links and an applicability inventory.
PDF extraction quality must be assessed on the actual source layout, especially
footnotes and tables.

MathDevMCP is useful after a property has been formalized. Its inspected Lean
route binds checks to an exact target and assumptions, and its result handling
distinguishes diagnostic evidence from proof. Its source-bound obligations and
resumable evidence structures can support the question in
Equation~\eqref{eq:refinement}. The legal adapter must define the rule language's
semantics, numeric and date operations, and the exact theorem requested. Structural
agreement between a formula and a code fragment is not a substitute for semantic
equivalence. The inspected project describes an internal research status and
has its own licensing terms; reuse within a bank needs the appropriate software
and ownership review.

% BEGIN REVISION ADDITION teach-P-04-prototype-03
\begin{figure}[H]
\centering
\TeachingCompare{ResearchAssistant: source processing}{MathDevMCP: scoped mathematics}{DynareMCP: bounded investigations}
\caption{Adjacent tools contribute different capabilities, each requiring local validation.}\label{fig:teach-P-04-prototype-03}
\end{figure}

% END REVISION ADDITION teach-P-04-prototype-03
DynareMCP supplies ideas for investigative planning, source graphs and meaningful
differences between model versions. Its economic-model parser and numerical
solvers are not direct engines for SFC requirements. The reusable pattern is a
reviewable change supported by a dependency graph, with unresolved alternatives
preserved. A legal adapter would compare rules, definitions, scope and dates
rather than macroeconomic equations.

The proposed search use is equally specific. Historical DynareMCP investigation
records use UCB-style exploration to choose work; the currently inspected
MathDevMCP derivation controller explicitly describes a different bounded
evidence-ordering method, not MCTS. It would be inaccurate to label both as the
same proof-search engine \citep{localrepos}.

For LegalMath, a search node could contain a candidate interpretation and its
unresolved questions. Actions might retrieve a referenced FAQ, request a
counterexample, check a definition, or prepare a question for a compliance owner.
A search score may prioritize expected information gain or reduction in review
work. It cannot determine which interpretation is legally correct. A branch with
more permissive outcomes or more successful solver calls receives no special
authority. If several readings remain viable, the product retains the disagreement.

Start with a bounded work queue ordered by materiality and dependencies. MCTS or
UCB scheduling becomes a candidate only if the prototype reveals a substantial
search-allocation bottleneck. Its evaluation would compare useful questions
resolved per unit of review effort under the same budget, while correctness
continues to depend on source review and formal evidence. Transaction execution
uses approved rules, with no exploratory search in the decision path.
% END SOURCE UNIT P-04-prototype:03



% BEGIN SOURCE UNIT P-05-evidence:02
\subsection{A reproducible research collection}
\label{merge:P-05-evidence:02}


The papers are stored under \path{docs/papers/} using the requested filename
convention: the title, an underscore, the first author's surname and the year in
parentheses. Filesystem-unsafe title punctuation is normalized, while the manifest
retains the original title. Each
record includes its source URL, bibliographic identity, edition note, retrieval
date, byte count, PDF page count and SHA-256 digest. The library index distinguishes
distinct works from alternate editions and identifies discovery leads that were
not available as usable full text.

The technical notes identify which part of each paper supports a proposed
borrowing and what remains untested. They distinguish a semantics from its
implementation, a reported experiment from a reproduced result, and a publicly
available repository from a licensed and approved bank dependency. Full-text
extraction is retained for local search, but the source PDF remains the reference
where layout or equations matter. ResearchAssistant extraction is an aid to
inspection, not certification of the extracted text.

% BEGIN REVISION ADDITION teach-P-05-evidence-02
\begin{figure}[H]
\centering
\TeachingFlow{Named local document copy}{Exact passage and source hash}{Manuscript claim and support judgment}
\caption{An audit trail must show what was actually read for a citation.}\label{fig:teach-P-05-evidence-02}
\end{figure}

% END REVISION ADDITION teach-P-05-evidence-02
The SFC source snapshots and the earlier local-repository inspection records
remain under \path{.localresources/}, with their own manifest. The LaTeX and
bibliography are now canonical in \path{docs/monograph/}; the earlier
proposal entry point builds the same unified work. This arrangement permits another
reviewer to trace the proposal's substantive examples back to a document edition
without reconstructing the web searches. Redistributing the paper library is a
separate matter from using the local copies for this research; the source licenses
and access terms remain relevant.
% END SOURCE UNIT P-05-evidence:02

\section{Exact controller sequence}
\label{sec:controller-reference}

\begin{lstlisting}[caption={Main controller, with explicit terminal uncertainty.}]
start_run(source_packet, profile, required_policy):
    validate_source_identity_and_policy()
    commit(run = INITIALIZING, counters = 0)
    reserve_and_issue_initial_roles_with_global_budget()
    collect_until_complete_or_action_deadline()
    validate_outputs_and_freeze_initial_proposals()
    reconcile_inventory_candidates_and_required_roles()

    while has_unresolved_issue():
        if integrity_failure():
            close_as(FAILED_INTEGRITY); break
        if cancelled():
            close_as(CANCELLED); break
        if deadline_or_global_or_round_limit_reached():
            mark_remaining_issues_with_stop_reason()
            close_as(BLOCKED_UNRESOLVED); break
        if no_progress_limit_reached():
            mark_remaining_issues(NO_PROGRESS)
            close_as(BLOCKED_UNRESOLVED); break

        plan = choose_permitted_discriminating_actions()
        if plan.is_empty():
            mark_remaining_issues(HUMAN_REQUIRED)
            close_as(BLOCKED_UNRESOLVED); break
        commit_next_round_and_reserve_bounded_actions(plan)
        dispatch_only_reserved_actions()
        collect_results_until_each_action_is_terminal()
        validate_results_and_create_candidate_revisions()
        rerun_affected_checks_and_reconcile_all_open_issues()
        record_progress_signature()

    if not already_closed():
        apply_integrity_cancellation_and_deadline_guards()
    if not already_closed():
        check_required_roles_coverage_and_all_mandatory_checks()
        if any_requirement_incomplete():
            close_as(BLOCKED_UNRESOLVED)
        else:
            close_as(READY_FOR_REVIEW)
    write_complete_report_for_every_terminal_status()
    return report_reference
\end{lstlisting}

\section{Machine-readable blocked-result example}

\begin{lstlisting}[caption={A blocked result preserves alternatives and the next decision.}]
{
  "schema_version": "interpretation.v1",
  "run_status": "BLOCKED_UNRESOLVED",
  "source_ref": "23EC46-retained-version",
  "selected_slice": "paragraph-10",
  "candidate_ids": ["cashback-reimbursement", "cashback-fee-reduction"],
  "material_unresolved_issue_ids": ["classification-cashback"],
  "release_eligible": false,
  "processing_stop": "ROUND_LIMIT",
  "rounds_issued": 3,
  "probability_of_legal_correctness": null,
  "next_decision": "Establish the applicable fee-discount definition",
  "next_decision_owner_role": "LEGAL_INTERPRETATION_REVIEWER"
}
\end{lstlisting}

\section{Exact financial-rule serialization}
\label{sec:financial-serialization}

The complete financial-condition bundle is supplied in
\path{fixtures/spi-financial.bundle.json} beneath the specification directory.
It declares two money inputs, an unconditional scope for this isolated
subcondition, the expression, interpretation record and a verified span in
Annex 1 paragraph 3.1. The excerpt below is a complete expression node from
that bundle. Node identifiers make its trace reproducible.

\begin{minipage}{\textwidth}
\begin{lstlisting}[caption={The portfolio comparison in the supplied JSON bundle. Amounts are HK cents.}]
{"node_id":"portfolio.ge", "op":"compare", "cmp":"ge",
 "left":{"node_id":"portfolio.input", "op":"fact",
         "name":"portfolio"},
 "right":{"node_id":"portfolio.threshold", "op":"literal",
          "type":"money_hkd", "value":"4000000000"}}
\end{lstlisting}
\end{minipage}

The second comparison uses \texttt{net\_assets\_ex\_home} and
\texttt{"8000000000"}; an \texttt{any} node combines them. These inputs denote
amounts normalized under adopted wealth definitions, not any account balance
supplied under the same name. Release requires those definitions and applicability
to be reviewed. The narrow result name, \texttt{spi.financial}, identifies a
subcondition rather than a complete SPI assessment or trade authorization.


\section{Complete rule-language grammar}

An implementer can express the grammar below directly as tagged records. Every
\texttt{Expr} additionally carries a unique \texttt{node\_id}. A rule is
\texttt{(id, type, scope, body, interpretation\_id, source\_span\_ids)};
\texttt{scope} is Boolean and \texttt{body} has the declared result type.
\begin{lstlisting}[caption={Complete operator grammar for RuleIR 0.1; repeated arguments are ordered. The Java demonstration implements a stated subset.}]
Type = bool | integer | money_hkd | date
Expr = literal(Type, value)
     | fact(name) | rule(name)
     | all([Expr, ...]) | any([Expr, ...]) | not(Expr)
     | compare(eq | ge | gt, left, right)
     | add(left, right) | sub(left, right)
     | scale(arg, numerator, denominator)
     | if(condition, then_value, else_value)
     | default(base, [Exception, ...])
Exception = (exception_id, guard, value,
             interpretation_id, source_span_ids)
\end{lstlisting}


\section{Durable request accounting and worker recovery}

An action moves through \texttt{RESERVED}, \texttt{DISPATCHED}, and one of
\texttt{SUCCEEDED}, \texttt{FAILED}, \texttt{TIMED\_OUT}, or
\texttt{RESULT\_UNKNOWN}. The last state is necessary when the controller cannot
determine whether a remote service performed the action. A lost connection after
submission is different from a request rejected before dispatch. Both can lead
to missing results, but only the former may already have consumed remote work.

The reservation transaction records the action identifier, input hashes, issue
root, round, budget charge, adapter identity, timeout and idempotency key. An
outbox entry is written in the same transaction. A worker claims that entry with
a lease and a fencing token. The token increases when ownership changes, so an
old worker cannot commit a result after a newer owner has taken responsibility.
A valid late result can be retained as an observation without changing the closed
action's authoritative state.

The dispatch boundary requires care. If the outbox records only ``not yet sent''
and a worker crashes after sending but before updating the database, another
worker may send again. A provider idempotency key can make that retry refer to the
same operation, but not every provider supports this guarantee. The adapter must
state its capability. Without remote idempotency, a recovery policy can choose to
leave the result unknown rather than risk another paid request, or authorize a
new counted action. It cannot truthfully promise exactly once remote execution.

% BEGIN REVISION ADDITION teach-M06-17
% Teaching figure remains in the main chapter.
% END REVISION ADDITION teach-M06-17
The budget system reserves an upper bound on tokens or cost when those resources
are capped. Settlement records actual usage when known. A refund of unused
monetary reservation may be legitimate after a confirmed terminal response, but
it does not decrement the issued-action count. Otherwise a sequence of failed
requests could evade the maximum number of attempts. If actual usage exceeds a
provider's documented bound, the run records an accounting discrepancy and stops
new spending until the integration is investigated.

A compare-and-swap version protects issue updates. The worker returns evidence
against the issue revision it received. If the issue has changed, the controller
revalidates whether the evidence is still relevant before applying a resolution.
This prevents two simultaneous actions from overwriting each other's objections
or resolving an issue using a candidate that is no longer current.

Cancellation stops new reservations and attempts to cancel outstanding work.
Remote cancellation may be best effort. The report records outstanding or late
actions and keeps their budget charges. Cancellation cannot convert an incomplete
run into a clean result, and it cannot delete the source of a later adverse
finding. These details matter because regulatory evidence must survive ordinary
operational events, including a person closing a browser or restarting a worker.


% BEGIN SOURCE UNIT P-02-literature:01
\subsection{Scope and method of this review}
\label{merge:P-02-literature:01}


The review began with the seven requested links and expanded through their
references, forward-citation searches for Catala and Stipula, and searches on
executable law, legal reasoning, formalization, dates, testing and process
compliance. ResearchAssistant supplied discovery records and local PDF extraction.
The saved OpenAlex citation searches returned 52 entries for the selected Catala
record and 13 for the selected Stipula journal record. Those are retrieval counts
for particular indexed editions, not complete or comparable measures of influence.
Preprints, revised editions and new publications can be missing or duplicated.

The initial proposal retained 38 paper editions representing 37 distinct
works. The first unified library retained 50 editions representing 49 works,
including the later legal-search and uncertainty studies. The present
audit adds two legal-reliability studies, bringing the academic collection
to 52 editions representing 51 works. Regulatory documents, standards
and software records are retained alongside them.
The AAAI tax-reasoning paper and its extended preprint are retained separately;
all seven requested sources are present. Technical reading covered the relevant
methods, semantics, evaluation and supporting appendix material. The accompanying
notes identify the inspected sections, version differences, implementation checks
and limits of adoption \citep{libraryreview}. This is a focused technical review
supporting a product decision, rather than a claim to have enumerated every
citation. Unavailable secondary leads are recorded separately and do not support
implementation claims.

% BEGIN REVISION ADDITION teach-P-02-literature-01
% Teaching figure remains with the main discussion.
% END REVISION ADDITION teach-P-02-literature-01
The expanded search treats eight families separately: executable statutes;
defaults and defeasible priorities; obligations and violations; temporal process
compliance; authority and provenance; legal languages, compilers and testing;
neurosymbolic formalization and evaluation; and deployable rule engines. Important
foundations predate Catala, while contract languages have partially separate
citation communities. The retained forward-citation records have individual
dispositions in \path{docs/papers/coverage.md}; additional searches and backward
references supply the missing mechanisms. Selection favors primary sources that
define a needed mechanism or report an informative failure. Downloading a paper
or screening its abstract is distinguished from inspecting its technical method.
This is a bounded review of the prototype's design questions. Remaining retrieval
gaps are stated rather than counted as read papers.

Version control is particularly important here. The ARc paper was first posted
in 2025, but the inspected second version is from July 2026. The eFLINT paper's
inspected version is from August 2026. The Stipula preprint requested by the
sponsor has three authors; the later journal edition has a different author list.
Bibliographic records and local filenames make these distinctions explicit.
Published results are reported as their authors' findings; none of their
benchmarks or proof developments was independently rerun for this proposal.
% END SOURCE UNIT P-02-literature:01

\section{Release identity and reproducible builds}
\label{sec:release-build-record}

A build manifest records the bundle, generated source, runtime source, toolchain,
Java target, dependencies and JAR digest. A verification report then names the
build and exact JAR it tested. A release manifest names that build and verification
report, together with applicability and effective-time information. Approval
binds the final release manifest and bundle.

The direction matters. Embedding a manifest that contains the JAR's own digest
inside that JAR creates a circular identity problem. The current design keeps
these records external and acyclic. A reviewer approves the exact tested bytes
through the final manifest, rather than approving a mutable directory or a
human-readable version label.

Candidate compilation occurs before legal approval, because reviewers need a
concrete package and verification evidence to inspect. Compilation does not grant
deployment authority. Export into a controlled release route occurs only after
the required approvals and checks. This distinction lets engineers do useful
reversible work without repeatedly asking for permission, while preserving the
human decision at the point where it matters.

Reproducible builds strengthen the identity chain. Under the same declared
toolchain and inputs, build in separate directories with fixed archive-entry
timestamps and compare emitted source and binary hashes. A mismatch should be
investigated and recorded. Reproducibility does not prove semantic correctness,
but it reduces uncertainty about which bytes correspond to the reviewed source.

The local MVP uses synthetic identities and unsigned historical archives to
exercise workflow. A bank deployment needs enterprise identity, role separation,
key management and independently verified release signatures or equivalent
controls. A hash proves equality of bytes under its cryptographic assumption;
it does not prove who authorized those bytes. The manuscript's implementation
plan keeps that integration distinct from the local demonstration.


\backmatter
\begingroup
\raggedright
\bibliographystyle{plainnat}
\bibliography{\LegalMathBibliography}
\endgroup
\end{document}
